\pdfoutput=1
\documentclass[10pt,twocolumn,letterpaper]{article}

\usepackage{wacv}
\usepackage{times}
\usepackage{epsfig}
\usepackage{graphicx}
\usepackage{amsmath}
\usepackage{amssymb}

% Include other packages here, before hyperref.
\usepackage{caption}
\usepackage{subcaption}
\usepackage{cleveref}
\usepackage{multirow}
\usepackage{booktabs,tabularx}
\newcolumntype{Y}{>{\raggedright\arraybackslash}X} 
\usepackage{tabularx}
\usepackage{enumerate}
\usepackage{url}

% If you comment hyperref and then uncomment it, you should delete
% egpaper.aux before re-running latex.  (Or just hit 'q' on the first latex
% run, let it finish, and you should be clear).
%\usepackage[pagebackref=true,breaklinks=true,letterpaper=true,colorlinks,bookmarks=false]{hyperref}

\wacvfinalcopy % *** Uncomment this line for the final submission

\def\wacvPaperID{601} % *** Enter the wacv Paper ID here
\def\httilde{\mbox{\tt\raisebox{-.5ex}{\symbol{126}}}}

% Pages are numbered in submission mode, and unnumbered in camera-ready
\ifwacvfinal\pagestyle{empty}\fi
\setcounter{page}{1}
\begin{document}
\newcommand{\MK}[1]{\textcolor{black}{#1}}
\newcommand{\PN}[1]{\textcolor{black}{#1}}

%%%%%%%%% TITLE
\title{PlotQA: Reasoning over Scientific Plots}

% Authors at the same institution
\author{Nitesh Methani \thanks{The first two authors have contributed equally} \hspace{1cm} Pritha Ganguly\textsuperscript{*} \hspace{1cm} Mitesh M. Khapra \hspace{1cm} Pratyush Kumar \\\\
Department of Computer Science and Engineering\\
Robert Bosch Centre for Data Science and AI (RBC-DSAI)\\
Indian Institute of Technology Madras, Chennai, India\\
{\tt\small \{nmethani, prithag, miteshk, pratyush\}@cse.iitm.ac.in}
}
% Authors at different institutions
% \author{First Author \\
% Institution1\\
% {\tt\small firstauthor@i1.org}
% \and
% Second Author \\
% Institution2\\
% {\tt\small secondauthor@i2.org}
% }

\newcommand{\revision}[1][\textcolor{black}]{#1}%modifications done for revision
\newcommand{\delete}[1][\textcolor{b}]{#1}%to be deleted as per revision

\maketitle
\ifwacvfinal\thispagestyle{empty}\fi

%%%%%%%%% ABSTRACT
\begin{abstract}
%Reasoning over plots by question answering (QA) is a challenging machine learning task at the intersection of vision, language processing, and reasoning.
Existing synthetic datasets (FigureQA, DVQA) for reasoning over plots do not contain variability in data labels, real-valued data, or complex reasoning questions. 
Consequently, proposed models for these datasets do not fully address the challenge of reasoning over plots. \MK{In particular, they assume that the answer comes either from a small fixed size vocabulary or from a bounding box within the image. However, in practice, this is an unrealistic assumption because many questions require reasoning and thus have real-valued answers which appear neither in a small fixed size vocabulary nor in the image. In this work, we aim to bridge this gap between existing datasets and real-world plots}. %by proposing a new dataset where a large fraction of answers are out-of-vocabulary words.} 
Specifically, we propose PlotQA with \revision{28.9 million question-answer pairs} over \textcolor{black}{224,377 plots} on data from real-world sources and questions based on crowd-sourced question templates. 
Further, \revision{80.76\%} of the out-of-vocabulary (OOV) questions in PlotQA have answers that are not in a fixed vocabulary.
\MK{Analysis of existing models on PlotQA reveals that they cannot deal with OOV questions:  their overall accuracy on our dataset is in single digits. This is not surprising given that these models were not designed for such questions. As a step towards a more holistic model which can address fixed vocabulary as well as OOV questions, we propose a hybrid approach}: Specific questions are answered by choosing the answer from a fixed vocabulary or by extracting it from a predicted bounding box in the plot, while other questions are answered with a table question-answering engine which is fed with a structured table generated by detecting visual elements from the image. 
%For the latter, we propose the VOES pipeline and combine it with SAN-VQA to form a hybrid model SAN-VOES.
On the existing DVQA dataset, our model has an accuracy of 58\%, significantly improving on the highest reported accuracy of 46\%. 
On PlotQA, our model has an accuracy of \revision{22.52\%}, which is significantly better than state of the art models. 
%Analysis of each module in the VOES pipeline reveals that further improvement in accuracy requires more accurate visual element detection.
\end{abstract}

%%%%%%%%% BODY TEXT
\section{Introduction}

Data plots such as bar charts, line graphs, scatter plots, etc. provide an efficient way of summarizing numerical information. 
% They are frequently used in textbooks, research papers, professional reports, and newspaper articles. 
%Interpreting and reasoning over such plots are considered tests of human aptitude. 
%It is thus of interest to formulate and evaluate machine comprehension of plots.
Recently, in \cite{FQA, DVQA} two datasets containing plots and deep neural models for question answering over the generated plots have been proposed. In both the datasets, the plots are synthetically generated with data values and labels drawn from a custom set. In the FigureQA dataset \cite{FQA}, all questions are binary wherein answers are either Yes or No, (see Figure \ref{fig:sample_plots_fqa} for an example).  The DVQA dataset \cite{DVQA}, generalizes this to include questions which can be answered either by (a) fixed vocabulary of 1000 words, or (b) extracting text (such as tick labels) from the plot.  An example question could seek the numeric value represented by a bar of a specific label in a bar plot (see Figure \ref{fig:sample_plots_dvqa}).  Given that all data values in the DVQA dataset are chosen to be integers and from a fixed range, the answer to this question can be extracted from the appropriate tick label. 

%This task lies in the intersection of vision, language processing, and reasoning, and thus poses interesting research challenges.
%It also has widespread real-world applicability. 
%A highly accurate model for plot reasoning can help domain experts such as policymakers and doctors access information in a large collection of plots.
%This can also help visually impaired persons interact with plots in natural language.
% Dataset comparison
\begin{figure*}
\centering
\begin{subfigure}[t]{.32\textwidth}
  \centering
  \includegraphics[width=.9\linewidth]{images/fqa.png}
\end{subfigure}%
\begin{subfigure}[t]{.32\textwidth}
  \centering
  \includegraphics[width=.9\linewidth]{images/dvqa.png}
\end{subfigure}
\begin{subfigure}[t]{.32\textwidth}
  \centering
  \includegraphics[width=1.02\linewidth]{images/plotqa.png}
\end{subfigure}
\begin{subfigure}[t]{.3\textwidth}
  \noindent
  \footnotesize
  Q: Is Light Green the minimum?\\
  A: $1$
  \caption{FigureQA}
  \label{fig:sample_plots_fqa}
\end{subfigure}
\begin{subfigure}[t]{.3\textwidth}
  \noindent 
  \footnotesize
  Q: What is the value of mad in drop?\\
  A: $7$
  \caption{DVQA}
  \label{fig:sample_plots_dvqa}
\end{subfigure}
\begin{subfigure}[t]{.3\textwidth}
  \noindent
  \footnotesize
  Q: What is the average number of Hispanic students in schools?
  A: $51.67$
  \caption{PlotQA}
  \label{fig:sample_plots_PlotQA}
\end{subfigure}
\caption{
A sample \{plot, question, answer\} triplet from FigureQA, DVQA, and PlotQA (our) datasets.
%In comparison to the FigureQA and DVQA, the triplets of the PlotQA dataset are more complex: plots are generated from real data, questions are more complex, and answers can be real-valued. 
}
\label{fig:sample_plots}
\end{figure*}

%----PlotQA : Question distribution and sample questions
\begin{table*}[ht!]
\begin{center}
\footnotesize
\begin{tabularx}{\textwidth}{|X|X|X|X|}
\hline
&\multicolumn{3}{|c|}{\textbf{Question Type}} \\ \cline{2-4}
\multicolumn{1}{|c|}{\textbf{Answer Type}} & \textbf{Structure} & \textbf{Data Retrieval} & \textbf{Reasoning} \\ \hline
\multirow{1}{*}{Yes/No}
 & \multicolumn{1}{l|}{\begin{tabular}[c]{@{}l@{}}Does the graph\\ contain grids?\end{tabular}} & \multicolumn{1}{l|}{\begin{tabular}[c]{@{}l@{}}Does the price of diesel in \\ Barbados monotonically\\ increase over the years?\end{tabular}} & \multicolumn{1}{l|}{\begin{tabular}[c]{@{}l@{}}Is the difference between the price of diesel in Angola in 2002 and 2004 \\greater than the difference between the price of diesel \\in Lebanon in 2002 and 2004?\end{tabular}} \\ \hline
\multirow{1}{*}{Fixed vocabulary}
 & \multicolumn{1}{l|}{\begin{tabular}[c]{@{}l@{}}How are the legend \\ labels stacked?\end{tabular}} & \multicolumn{1}{l|}{\begin{tabular}[c]{@{}l@{}}What is the label or title of \\ the X-axis ?\end{tabular}} & \multicolumn{1}{l|}{\begin{tabular}[c]{@{}l@{}}In how many years, is the price of diesel greater than 0.6 units?\end{tabular}} \\ \hline
\multirow{1}{*}{Open vocabulary}
 & - & \multicolumn{1}{l|}{\begin{tabular}[c]{@{}l@{}}What is the price of diesel in \\ Lebanon in the year 2008?\end{tabular}} & \multicolumn{1}{l|}{\begin{tabular}[c]{@{}l@{}}What is the ratio of the price of diesel in Lebanon in 2010 to that in 2014?\end{tabular}} \\ \hline
\end{tabularx}
\caption{\revision{Sample questions for 9 different question-answer types in PlotQA. The example questions are with respect to the plot in Figure \ref{fig:dataset_examples_line}. Note that there are no open vocabulary answers for \textit{Structural Understanding} questions.}}
\label{sample_qs_stats}
\end{center}
\end{table*}
%-------------------------------------------------------

%-------------------------------------------------------------------------
% DVQA vs FigureQA vs PlotQA
\begin{table*}[]
\footnotesize
\begin{center}
\begin{tabular}{|l|c|c|c|c|c|c|c|c|}
\hline
\multicolumn{1}{|c|}{\textbf{Datasets}} & \textbf{\begin{tabular}[c]{@{}c@{}}\#Plot \\ types\end{tabular}} & \textbf{\begin{tabular}[c]{@{}c@{}}\#Plot\\ images\end{tabular}} & \textbf{\begin{tabular}[c]{@{}c@{}}\#QA \\ pairs\end{tabular}} & \textbf{\begin{tabular}[c]{@{}c@{}}Vocabulary \end{tabular}} & \textbf{\begin{tabular}[c]{@{}c@{}}Avg. question \\ length\end{tabular}} & \textbf{\#Templates} & \textbf{\begin{tabular}[c]{@{}c@{}}\#Unique \\ answers\end{tabular}} & \textbf{\begin{tabular}[c]{@{}c@{}}Open \\ vocab.\end{tabular}} \\ \hline
FigureQA & 4 & 180,000 & 2,388,698 & \begin{tabular}[c]{@{}c@{}}100 colours from\\ X11 colour set\end{tabular} & 7.5 & \begin{tabular}[c]{@{}c@{}}15 \\ (no variations)\end{tabular} & 2 & Not present \\ \hline
DVQA & 1 & 300,000 & 3,487,194 & \begin{tabular}[c]{@{}c@{}}1K nouns from\\ Brown corpus\end{tabular} & 12.30 & \begin{tabular}[c]{@{}c@{}}26 \\ (without paraphrasing)\end{tabular} & 1576 & Not present \\ \hline
PlotQA & 3 & 224,377 & 28,952,641 & \begin{tabular}[c]{@{}c@{}}Real-world  axes variables \\ and floating point numbers \end{tabular} & 43.54  & \begin{tabular}[c]{@{}c@{}}74 \\ (with paraphrasing)\end{tabular} & 5,701,618 & Present \\ \hline
\end{tabular}
\caption{\revision{Comparison between the existing datasets (FigureQA and DVQA) and our proposed dataset (PlotQA). }}
\label{tab:comparison_data_stats}
\end{center}
\end{table*}
%----------------------------------------------

While these datasets have initiated the research questions on plot reasoning, realistic questions over plots are much more complex. For instance, consider the question in Figure \ref{fig:sample_plots_PlotQA}, where we are to compute the average of floating point numbers represented by three bars of a color specified by the label.
The answer to this question is neither in a fixed vocabulary nor can it be extracted from the plot itself. 
Answering such questions requires a combination of perception, language understanding, and reasoning, and thus poses a significant challenge to existing systems.
Furthermore, this task is harder if the training set is not synthetic, but instead is sourced from real-world data with large variability in floating-point values, large diversity in axis and tick labels, and natural complexity in question templates.

To address this gap between existing datasets and real-world plots, we introduce the PlotQA\footnote{Dataset can be downloaded from \url{bit.ly/PlotQA}} dataset with \revision{28.9 million} question-answer pairs grounded over \textcolor{black}{224,377} plots.
PlotQA improves on existing datasets on three fronts.
First, roughly \revision{80.76\%} of the questions have answers which are not present in the plot or in a fixed vocabulary. 
Second, the plots are generated from data sourced from World Bank, government sites, etc., thereby having a large vocabulary of axis and tick labels, and a wide range in data values. 
Third, the questions are complex as they are generated based on 74 templates extracted from 7,000 crowd-sourced questions asked by workers on a sampled set of 1,400 plots. 
Questions are categorized into 9 (=3x3) cells based on the question type: `Structural Understanding', `Data Retrieval', or `Reasoning' and and the answer type: `Yes/No', `From Fixed Vocabulary', or `Out Of Vocabulary (OOV)' (see Table \ref{sample_qs_stats}). 
%\revision{Sample questions and the fraction of questions for these cells are shown in Table \ref{sample_qs_stats} and Table \ref{template_ans_stats} respectively.}

\MK{We first evaluate three state of the art models on PlotQA, \textit{viz.}, SAN-VQA\cite{SAN}, Bilinear attention network (BAN) \cite{DBLP:conf/nips/KimJZ18} and LoRRA \cite{singh2019TowardsVM}. Note that, by design none of these models are capable of answering OOV questions.
% but we still evaluate them for the sake of completeness and for empirically establishing the need for better models. 
In particular, SAN-VQA and BAN treat plot reasoning as a classification task and expect the answer to lie in a small vocabulary whereas in our dataset the answer vocabulary is prohibitively large ($\sim$5M words). Similarly, LoRRA assumes that the answer is present in the image itself as a text and the task is to just extract this region containing the text followed by OCR (optical character recognition). Again, such a model will be unable to answer questions such as the one shown in Figure \ref{fig:sample_plots_PlotQA}, which form a significant segment of real-world use-cases and our dataset. As a result, these these models give an accuracy of less than 8\% on our dataset. On the other hand, existing models (in particular, SAN) perform well on questions with answers from a fixed vocabulary, %including questions with Yes/No answers, 
which was the intended purpose of these models.}

%First, SAN-VQA \cite{SAN} performs well on Structural Understanding questions and on Data Retrieval questions if the answer is binary.
%Second, SAN-VQA performs poorly for Reasoning questions. In particular, it is unable to answer any reasoning question where the answer is from an open vocabulary.

\MK{Based on the above observations, we propose a hybrid model with a binary classifier which given a question decides if the answer would lie in a small top-$k$ vocabulary or if the answer is OOV. For the former, the question is passed through a classification pipeline which predicts a distribution over the top-$k$ vocabulary and selects the most probable answer. For the latter (arguably harder questions), we pass the question through a pipeline of four modules: Visual element detection, Optical character recognition, Extraction into a structured table, and Structured table question answering.}
%For questions with an answer from an open vocabulary, VOES correctly answers \revision{57.39\%} of Data Retrieval questions and \revision{14.95\%} of the Reasoning questions. 
%Given the complementary strengths of SAN-VQA and VOES, we train a hybrid model with a binary classifier which given a question decides whether to use the SAN-VQA or the VOES model. 
This proposed hybrid model significantly outperforms the existing models and has an aggregate accuracy of \revision{22.52\%} on the PlotQA dataset.
We also evaluate our model on the DVQA dataset where it gives an accuracy of 58\%, improving on the best-reported result of SANDY \cite{DVQA} of 46\%.
% On the PlotQA dataset, we analyze the performance of VOES-Oracle wherein the ground-truth structured table is available to the question-answering model.
%\delete{On the PlotQA dataset, we analyze the performance of VOES-Oracle wherein we feed the ground truth structured table to the question-answering model. This ground truth structured table is constructed from the ground truth annotations of the plots instead of using the annotations predicted by our perception model.
%VOES-Oracle performs significantly better than VOES, highlighting the need to improve the perception module to increase end-to-end accuracy.
% \textcolor{black}{We believe that this engineering effort is necessary to establish a baseline for this dataset and highlight the above lacunae in existing SOTA models for object detection.}
In summary, we make two major contributions:\\
(1) We propose PlotQA dataset with plots on data sourced from the real-world and questions based on templates sourced from manually curated questions.
The dataset exposes the need to train models for questions that have answers from an Open Vocabulary. \\
(2) We propose a hybrid model with perception and QA modules for questions that have answers from an Open Vocabulary. This model gives the best performance not only on our dataset but also on the existing DVQA dataset.
%    VOES performs significantly better on these questions than existing models. \\
%(3) We propose a hybrid model SAN-VOES that combines the strength of classification and extraction methods (SAN) with VOES. 
%    SAN-VOES significantly improves on the best results on both DVQA and PlotQA. \\
%\textcolor{black}{
%(4) We empirically show that detecting visual elements from plot images is still an open challenge and generating structured tables from plots is a difficult CV task.}
%\end{itemize}

% \begin{itemize}
%     \item \textcolor{blue}{Motivation for introducing PlotQA: What are the already existing datasets; what do they address and what do they lack?} \\
%     While doing question answering on scientific plots, the answer can either :
%     \begin{itemize}
%         \item be binary, i,e Yes or No
%         \item come from a set of fixed vocabulary.
%         \item has to be calculated, extracted or generated from the plot's data, for example, average of numbers, difference of numbers, axis labels of the plots \textit{etc}
%     \end{itemize}
%     Existing datasets \textit{viz.}, FigureQA \cite{FQA} and DVQA  \cite{DVQA} have only those questions types for which answers can either be classified from a fixed vocabulary or are plot specific,i.e the text can be extracted from the plot. The FigureQA \cite{FQA} dataset has questions where the answers can be only Yes or No. The DVQA \cite{DVQA} dataset generalizes this to include questions which can be answered either from a vocabulary of about 1000 words or by extracting text (such as tick labels) from the plot, which makes the answers to the questions be classified from a fixed vocabulary of $1076$ words. However, consider the question given in Figure \ref{fig:sample_plots_PlotQA}. The answer to this question would neither be present in a fixed vocabulary nor can be extracted from text within the plot. Instead it needs to be computed by reasoning on different plot elements. FigureQA and DVQA do not have such reason-and-compute questions with out-of-vocabulary answers.
%     \item \textcolor{blue}{To overcome the above above shortcomings, we introduce PlotQA. What are the salient features of PlotQA? :}
%     \begin{itemize}
%         \item The data comes from real-world. Crowd-sourced templated, more complex questions. Out of $74$ templates,the answers to $23$ templates needs to calculated or generated from the plot's data.
%         \item In addition, we also include a novel test-set which contains plots based on data extracted from Open Government Data. This test set contains additional variable and legend names which are not seen in the World Bank Data and thus can check for model generalisability.
%     \end{itemize}
    
% \item \textcolor{blue}{Motivation for a new baseline, How do existing methods solve this problem? What are their shortcomings? :}
% Existing methods address only those questions types for which answers can either be classified from a fixed vocabulary or can be extracted text from the plot. For example, the Stacked Attention Network (SAN) \cite{SAN} classfies the answer from a fixed vocabulary. The models SANDY, MOM, proposed in the paper \cite{DVQA} classifies the answer from a fixed vocabulary and extracts the text from the plot. These existing models are incapable of answering questions whose answers need to calculated or generated.

% \item \textcolor{blue}{Our proposal of a pipeline based solution for the same: brief explanation of each module of the pipeline. Shortcomings of the pipeline. How do we overcome the shortcomings by introducing a Hybrid model.}

% \item \textcolor{blue}{Our overall contributions:}
% \begin{enumerate}
%         \item to the best of our knowledge,we are the first one to show that existing datasets do not address open vocabulary QAs
%         \item existing E2E VQA methods fails to do reasoning over OOV answers
%         \item introduced a dataset which has open vocabulary
%         \item proposed a SOTA baseline for plot QA
%     \end{enumerate}
% \end{itemize}
%------------------------------------------------------------------------

% Sample plots of PlotQA
\begin{figure*}
\centering
\begin{subfigure}[t]{.3\textwidth}
  \centering
  \includegraphics[width=.98\linewidth]{images/8_hbar_categorical.png}
  \caption{Horizontal bar graph}
  \label{fig:dataset_examples_hbar}
\end{subfigure}
\begin{subfigure}[t]{.3\textwidth}
  \centering
  \includegraphics[width=.95\linewidth]{images/line_sample.png}
  \caption{Line plot}
  \label{fig:dataset_examples_line}
\end{subfigure}
\begin{subfigure}[t]{.3\textwidth}
  \centering
  \includegraphics[width=1.05\linewidth]{images/dot_line_sample.png}
  \caption{Dot-Line graph}
  \label{fig:dataset_examples_dotline}
\end{subfigure}
\caption{Sample plots of different types in the PlotQA dataset.}
\label{fig:dataset_examples}
\end{figure*}
%-------------------------------------------------------


\section{Related Work}
%In this section, we review existing datasets and models for Visual QA.

\noindent\textbf{Datasets:} 
Over the past few years several large scale datasets for Visual Question Answering have been released. 
These include datasets such as COCO-QA \cite{RenKZ15}, DAQUAR \cite{MalinowskiF14}, VQA \cite{AntolALMBZP15,GoyalKSBP17} which contain questions asked over natural images. 
On the other hand, datasets such as CLEVR \cite{JohnsonHMFZG17} and NVLR \cite{SuhrLYA17} contain complex reasoning based questions on synthetic images having 2D and 3D geometric objects.  
There are some datasets \cite{KembhaviSKSHF16,KembhaviSSCFH17} which contain questions asked over diagrams found in text books but these datasets are smaller and contain multiple-choice questions. 
FigureSeer \cite{SiegelHLDF16} is another dataset which contains images extracted from research papers but this is also a relatively small (60,000 images) dataset. 
Further, FigureSeer focuses on answering questions based on line plots as opposed to other types of plots such as bar charts, scatter plots, \textit{etc.} as seen in %two recently released datasets, \textit{viz.}, 
FigureQA \cite{FQA} and DVQA  \cite{DVQA}. \MK{There is also the recent TextVQA \cite{singh2019TowardsVM} dataset which contains questions which require models to read the text present in natural images. This dataset does not contain questions requiring numeric reasoning. Further, the answer is contained as a text in the image itself.}
Thus, \emph{no} existing dataset contains plot images with complex questions which require reasoning and have answers from an Open Vocabulary.
%As discussed in the earlier section, these two datasets impose constraints which we have tried to generalize in the proposed DIP dataset. To the best of our knowledge, DIP dataset is the first work which introduces plots created from real-world data, containing natural $x$-$y$ axes variables, real data values, a large answer vocabulary and richer questions derived from crowdsourced instances.

\noindent \textbf{Models:}
The availability of the above mentioned datasets has facilitated the development of complex end-to-end neural network based models (\cite{SAN}, \cite{LuYBP16}, \cite{YangHGDS16}, \cite{NohH16}, \cite{SantoroRBMPBL17}, \revision{\cite{DVQA}, \cite{singh2019TowardsVM}}).
These end-to-end networks contain (a) encoders to compute a representation for the image and the question, (b) attention mechanisms to focus on important parts of the question and image, (c) interaction components to capture the interactions between the question and the image, \revision{(d) OCR module to extract the image specific text and (e) a classification layer for selecting the answer either from a fixed vocabulary or from a OCR appended vocabulary. }
% \revision{By design, these algorithms cannot generate answers which need to be computed and come from an open vocabulary.} \MK{Hence, they perform poorly on the PlotQA dataset as demonstrated empirically in section \ref{sec:experiments}.}% (as is the case for many questions in our dataset). %they fail to perform well on datasets where the questions involve more complex numeric reasoning over images. 
%A case in point is the performance of the models on the recently released datasets for plot QA, \textit{viz.}, FigureQA \cite{FQA} and DVQA  \cite{DVQA}.  
%One crucial difference is that when dealing with plots such as bar charts, scatter plots, etc. the numeric information in the image needs to be accurately inferred based on the position of the axes, relative position of legends, bars, lines etc. in the figure. 
%Such structure is typically not present in natural images and the resulting questions are also more complex and quantitative in nature. 
%As mentioned in the earlier section, an additional challenge in our dataset is that the answer vocabulary is not fixed and can range from floating point numerical answers to textual answers containing multiple words. 
%As a result, existing end-to-end models which treat VQA as a multi-class classification problem cannot be used. 
%-------------------------------------------------------------------------
\section{The PlotQA dataset}

In this section, we describe the PlotQA dataset and the process to build it. 
Specifically, we discuss the four main stages, \textit{viz.}, (i) curating data such as year-wise rainfall statistics, country-wise mortality rates, \textit{etc.}, (ii) creating different types of plots with a variation in the number of elements, legend positions, fonts, \textit{etc.}, (iii) crowd-sourcing to generate questions, and (iv) extracting templates from the crowd-sourced questions and instantiating these templates using appropriate phrasing suggested by human annotators.

\subsection{Data Collection and Curation}

We considered online data sources such as World Bank Open Data%\footnote{https://data.worldbank.org/}
, Open Government Data%\footnote{https://www.india.gov.in/}
, Global Terrorism Database%\footnote{https://www.start.umd.edu/gtd/}
%, Nutritional Analysis Data\footnote{https://www.kaggle.com/mcdonalds/nutrition-facts/home}
, \textit{etc.} which contain statistics about various indicator variables such as fertility rate, rainfall, coal production, \textit{etc.} across years, countries, districts, \textit{etc.} 
We crawled data from these sources to extract different variables whose relations could then be plotted (for example, rainfall v/s years across countries, or movie v/s budget, or carbohydrates v/s food\_item). 
%Some statistics about the crawled data are of interest. 
There are a total of 841 unique indicator variables (CO2 emission, Air Quality Index, Fertility Rate, Revenue generated, etc.) with 160 unique entities (cities, states, districts, countries, movies, food, etc.).
The data ranges from 1960 to 2016, though not all indicator variables have data items for all years.
The data contains positive integers, floating point values, percentages, and values on a linear scale. These values range from $0$ to $3.50\mathrm{e+}15$. 

%Data Source: Synthetic v/s WB, etc \\
%Axes labels: Randomly grouped v/s from real-world data \\
%Numeric values: Integers v/s Integers and Floating point numbers \\



\subsection{Plot Generation}
We included 3 different types of plots in this dataset, \textit{viz.}, bar plots, line plots, and scatter plots.
Within bar plots, we have grouped them by orientation as either horizontal or vertical. 
% Though not exhaustive, these plot types cover a major fraction of plots that we encounter in journals, research papers, and textbooks.
% Within the data sources we explored, we did not find enough data to create certain other types of plots such as Venn diagrams and pie charts which are used in specific settings. 
% We also do not consider composite plots such as Pareto charts which have line graphs on top of bar graphs. 
% Lastly, all the plots in our dataset contain only 2-axes. 
Figure \ref{fig:dataset_examples} shows one sample of each plot type. 
Each of these plot types can compactly represent 3-dimensional data. 
For instance, in Figure \ref{fig:dataset_examples_line}, the plot compares the \texttt{diesel prices} across \texttt{years} for different \texttt{countries}. 
To enable the development of supervised modules for various sub-tasks, we provide bounding box annotations for legend boxes, legend names, legend markers, axes titles, axes ticks, bars, lines, and title. 
By using different combinations of indicator variables and entities (years, countries, \textit{etc.}) we created a total of $224,377$ plots. 

To ensure variety in the plots, we randomly chose the following parameters: grid lines (present/absent), font size, notation used for tick labels (scientific-E notation or standard notation), line style (solid, dashed, dotted, dash-dot), marker styles for marking data points (asterisk, circle, diamond, square, triangle, inverted triangle), position of legends (bottom-left, bottom-centre, bottom-right, center-right, top-right), and colors for the lines and bars from a set of 73 colors. 
The number of discrete elements on the $x$-axis varies from 2 to 12 and the number of entries in the legend box varies from 1 to 4. 
%In other words, in the case of line plots, the number of lines varies from 1 to 4 and in the case of grouped bars the number of bars grouped on a single $x$-tick varies from 1 to 4. 
% For example, for the plots in Figure \ref{fig:dataset_examples_line}, the number of discrete elements on the $x$-axis is $6$ and the number of legend names (\textit{i.e.}, number of lines) is $3$.

%-------------------------------------------------------------------------------
% Added detailed stats for the extended paper on ArxiV
\begin{table*}[ht!]
\footnotesize
\begin{center}
\begin{tabular}{|l|rrrr|rrr|rrr|}
\hline
\multirow{2}{*}{\textbf{\begin{tabular}[c]{@{}l@{}}Dataset\\Split\end{tabular}}} & \multicolumn{4}{c|}{\textbf{Plot Types}}                                                                                                              & \multicolumn{3}{c|}{\textbf{Question Types}}                                                                                      & \multicolumn{3}{c|}{\textbf{Answer Types}}                                                                                    \\ \cline{2-11} 
& \multicolumn{1}{l}{\textbf{vbar}} & \multicolumn{1}{l}{\textbf{hbar}} & \multicolumn{1}{l}{\textbf{line}} & \multicolumn{1}{l|}{\textbf{dot-line}} & \multicolumn{1}{l}{\textbf{Structural}} & \multicolumn{1}{l}{\textbf{Data-Retrieval}} & \multicolumn{1}{l|}{\textbf{Reasoning}} & \multicolumn{1}{l}{\textbf{Yes/No}} & \multicolumn{1}{l}{\textbf{Fixed vocab.}} & \multicolumn{1}{l|}{\textbf{Open vocab.}} \\ \hline
Train      & 52,463 & 52,700 & 25,897 & 26,010 & 871,782 & 2,784,041 & 16,593,656 & 784,115 & 3,095,774 & 16,369,590 \\ %\hline
Validation & 11,249 & 11,292 & 5,547 & 5,571 & 186,994 & 599,573 & 3,574,081 & 167,871 & 600,424 & 3,592,353 \\ %\hline
Test       & 11,242 & 11,292 & 5,549 & 5,574 & 186,763 & 596,359 & 3,559,392 & 167,727 & 667,742 & 3,507,045 \\ \hline
\end{tabular}
\caption{\PN{Detailed Statistics for different splits of the PlotQA dataset.}}
\label{tab:detailed_stats}
\end{center}
\end{table*}

%-------------------------------------------------------------------------------


%-------------------------------------------------------
\begin{table}[]
\footnotesize
\begin{center}
\begin{tabular}{|l|r|r|r|}
\hline
&\multicolumn{3}{|c|}{\textbf{Question (Q) Type}} \\ \cline{2-4}
\multicolumn{1}{|c|}{\textbf{Answer (A) Type}} & \multicolumn{1}{c|}{\textbf{Structure}} & \multicolumn{1}{c|}{\textbf{Data Retrieval}} & \multicolumn{1}{c|}{\textbf{Reasoning}} \\ \hline
Yes/No & 36.99\% & 5.19\% & 2.05\% \\
Fixed  vocabulary & 63.01\% & 18.52\% & 15.92\% \\
Open vocabulary & 0.00\% & 76.29\% & 82.03\% \\ \hline
\end{tabular}
\caption{\PN{Overall distribution of Q and A types in PlotQA.}}
\label{template_ans_stats}
\end{center}
\end{table}
%-------------------------------------------------------

\subsection{Sample Question Collection by Crowd-sourcing}
As the source data of PlotQA dataset is significantly richer in comparison to FigureQA and DVQA, we found it necessary to ask a larger set of annotators to create questions over these plots. 
However, creating questions for all the plots in our dataset would have been prohibitively expensive. 
We sampled $1,400$ plots across different types and asked workers on Amazon Mechanical Turk to create questions for these plots. 
We showed each plot to 5 different workers resulting in a total of $7,000$ questions. 
We specifically instructed the workers to ask complex reasoning questions which involved reference to multiple plot elements in the plots. 
% We also gave the workers a list of simple questions such as ``Is this a vertical bar graph?'', ``What is the title of the graph?'', ``What is the value of coal production in 1993?'' and asked them to not create such questions as we had already created such questions using hand designed templates.  
We paid the workers USD $0.1$ for each question.

% %-------------------------------------------------------------------------
% \begin{table*}[ht!]
% \footnotesize
% \begin{center}
% \begin{tabular}{|l|c|c|c|c|c|c|c|c|}
% \hline
% \multicolumn{1}{|c|}{\textbf{Datasets}} & \textbf{\begin{tabular}[c]{@{}c@{}}\#Plot \\ types\end{tabular}} & \textbf{\begin{tabular}[c]{@{}c@{}}\#Plot\\ images\end{tabular}} & \textbf{\begin{tabular}[c]{@{}c@{}}\#QA \\ pairs\end{tabular}} & \textbf{\begin{tabular}[c]{@{}c@{}}Vocabulary \end{tabular}} & \textbf{\begin{tabular}[c]{@{}c@{}}Avg. question \\ length\end{tabular}} & \textbf{\#Templates} & \textbf{\begin{tabular}[c]{@{}c@{}}\#Unique \\ answers\end{tabular}} & \textbf{\begin{tabular}[c]{@{}c@{}}Open \\ vocab.\end{tabular}} \\ \hline
% FigureQA & 4 & 180,000 & 2,388,698 & \begin{tabular}[c]{@{}c@{}}100 colours from\\ X11 colour set\end{tabular} & 7.5 & \begin{tabular}[c]{@{}c@{}}15 \\ (no variations)\end{tabular} & 2 & Not present \\ \hline
% DVQA & 1 & 300,000 & 3,487,194 & \begin{tabular}[c]{@{}c@{}}1K nouns from\\ Brown corpus\end{tabular} & 12.30 & \begin{tabular}[c]{@{}c@{}}26 \\ (without paraphrasing)\end{tabular} & 1576 & Not present \\ \hline
% PlotQA & 3 & 224,377 & 28,952,641 & \begin{tabular}[c]{@{}c@{}}Real-world  axes variables \\ and floating point numbers \end{tabular} & 43.54  & \begin{tabular}[c]{@{}c@{}}74 \\ (with paraphrasing)\end{tabular} & 5,701,618 & Present \\ \hline
% \end{tabular}
% \caption{\revision{Comparison between the existing datasets (FigureQA and DVQA) and our proposed dataset (PlotQA). }}
% \label{tab:comparison_data_stats}
% \end{center}
% \end{table*}
% %----------------------------------------------


\subsection{Question Template Extraction \& Instantiation}

We manually analyzed the questions collected by crowdsourcing and divided them into a total of 74 templates. % (including the simple templates that we had manually designed as mentioned earlier).
These templates were divided into 3 question categories. 
These question categories along with a few sample templates are shown below. See Table \ref{template_ans_stats} for statistics of different question and answer types in our dataset (please refer to the Supplementary material for further details).

\noindent\textbf{Structural Understanding}: 
    These are questions about the overall structure of the plot and do not require any quantitative reasoning. 
    Examples: ``How many different coloured bars are there?'', ``How are the legend labels stacked?''.

\noindent\textbf{Data Retrieval}:
    These questions seek data item for a single element in the plot.
    Examples: ``What is the number of tax payers in Myanmar in 2015?''.%, ``How many bars are there on the $4^{th}$ tick from the top?''.

\noindent\textbf{Reasoning}: 
    These questions either require numeric reasoning over multiple plot elements or a comparative analysis of different elements of the plot, or a combination of both to answer the question. 
    Examples: 
    % ``Across all countries, what is the maximum price of diesel?'', 
    % ``In how many countries, is the cereal yield greater than 18000?'', 
    % ``What is the ratio of the value added by industrial sector in 1988 to that in 1989?''.
    ``In which country is the number of threatened bird species minimum?'', 
    ``What is the median banana production?'', 
    % ``What is the difference between the number of deaths in Bulgaria and Cuba in 2005?''.%, 
    %``In how many years, is the rice production greater than the average rice production over all years?''.

%For questions belonging to each of these 3  categories,  Table \ref{xx} shows the number of questions belonging to each of these categories. 
% For all questions, we created corresponding logical forms which could be executed on the raw data to extract the answers for the questions. 
% For example, for the question, ``Across all countries, what is the maximum number of threatened fish species?'', the corresponding logical form \cite{PasupatL15} would be, ``(max ((reverse fb:cell.cell.number) ((reverse fb:row.row.fish\_species) (fb:type.object.type fb:type.row))))''.

We abstracted the questions into templates such as ``In how many $<$plural form of X\_label$>$, is the $<$Y\_label$>$ of/in $<$legend\_label$>$ greater than the average $<$Y\_label$>$ of/in $<$legend\_label$>$ taken over all $<$plural form of X\_label$>$?''. We could then generate multiple questions for each template by replacing X\_label, Y\_label, legend\_label, etc. by indicator variables, years, cities etc. from our curated data. However, this was a tedious task requiring a lot of manual intervention. For example, consider the indicator variable ``Race of students'' in Figure \ref{fig:sample_plots_PlotQA}. 
If we substitute this indicator variable as it is in the above template, it would result in a question, ``In how many cities, is the race of the students(\%) of Asian greater than the average race of the students (\%) of Asian taken over all cities?'', which sounds unnatural. To avoid this, we asked in-house annotators to carefully paraphrase these indicator variables and question templates. %so that the variables could be automatically substituted in the template. 
The paraphrased version of the above example was ``In how many cities, is the percentage of Asian students greater than the average percentage of Asian students taken over all cities?''. Such paraphrasing for every question template and indicator variable required significant manual effort. Using this semi-automated process we generated a total of \revision{$28,952,641$} questions. %Given the diverse question templates and indicator variables, such paraphrasing was only possible with the help of in-house annotators.  %As a result, when each template is instantiated for different indicator variables it differs from its surface form.
This approach of creating questions on real-world plot data with carefully curated question templates followed by manual paraphrasing is a key contribution of our work. 
The resultant PlotQA dataset is much closer to the real-world challenge of reasoning over plots, significantly improving on existing datasets.
% \textcolor{black}{We believe that this approach of creating questions based on templates extracted from complex human generated questions is a good middle ground between  (i) the expensive and time consuming process of creating questions with the help of humans and (ii) the inexpensive and fast process of creating questions from very simple templates (as in FigureQA and DVQA).} 
\MK{Table \ref{tab:comparison_data_stats} summarizes the differences between PlotQA and these existing datasets such as FigureQA and DVQA. Note that (a) the number of unique answers in PlotQA is very large, (b) the questions in PlotQA are much longer, and (c) the vocabulary of PlotQA is more realistic than FigureQA or DVQA.}  
%Table \color[red]{\ref{,,}} shows the main differences between PlotQA and DVQA, highlighting the additional complexity introduced in our dataset.  

%\revision{As shown in Table \ref{fig:que_len}, the length of a question in the PlotQA dataset can vary from a minimum length of $6$ words to a maximum length of $84$ words, whereas FigureQA and DVQA datasets have questions with a maximum length of $11$ and $23$ respectively. Table \ref{template_ans_stats} shows that the answer to a question in the PlotQA dataset could either be Yes/No or from Fixed Vocabulary, or Open Vocabulary.}


\section{Proposed Model}
\label{sec:proposed_model}
%\subsection{VQA Baseline}
%As a baseline, we run the state of the art VQA algorithm, Stacked Attention Network (SAN) \cite{SAN} on PlotQA. The SAN uses a multi-layer attention \cite{Attention} mechanism that uses the semantic representation of the question to query the plot image and locate relevant visual regions to infer the correct answer. The SAN network is also an encoder-decoder framework where the image features are extracted from the last pooling layer of a deep Convolutional Neural Network, which is further combined with the LSTM \cite{LSTM} features of the question to form a query vector. The query vector is used to locate relevant visual regions which is used to compute a softmax over a fixed vocabulary of answers.

%\subsection{Our Proposed Pipeline}
%\textcolor{red}{This subsections has been copied from CVPR paper.}\\
\MK{Existing models for VQA are of two types: (i) read the answer from the image (as in LoRRA) or (ii) pick the answer from a fixed vocabulary (as in SAN and BAN).} Such models work well for datasets such as DVQA where indeed all answers come from a fixed vocabulary (global or plot specific) but are not suited for PlotQA with a large number of OOV questions. %However, in our dataset, for roughly \revision{76.29\%} of \textit{Data Retrieval} questions and \revision{82.03\%} of \textit{Reasoning} questions, the answers do not come from a fixed vocabulary, but instead need to be computed by reasoning over one or more visual elements in the plots.
\MK{Answering such questions involves various sub-tasks: (i) detect all the elements in the plot (bars, legend names, tick labels, etc), (ii) read the values of these elements, (iii) establish relations between the plot elements, e.g., creating tuples of the form \{country=Angola, year=2006, price of diesel = 0.4 \}, and (iv) reason over this structured data. Expecting a single end-to-end model to be able to do all of this is unreasonable.}
%To address such complex questions, we seek to leverage existing results on QA over tables \cite{SEMPRE}. 
%However, this requires the intermediate step of translating the plot image into a structured table, potentially similar to the one from which the plot was generated.
Hence, we propose a multi-staged pipeline to address each of the sub-tasks.
% To this end, we propose a method which separates visual element detection from reasoning for QA. 
% Our multi-stage approach contains modules for (a) detecting visual elements in the plot such as bars, bounding boxes around axes label, etc. (b) performing optical character recognition within these bounding boxes, (c) converting this data into a structured table, and (d) answering questions using this structured table. %In the following, we describe each of these modules in detail.

\MK{We further note that for simpler questions which do not require reasoning and can be answered from a small fixed size vocabulary, such an elaborate pipeline is an overkill. As an illustration consider the question ``How many bars are there in the image?''. This does not require reasoning and can be answered based on visual properties of the image. For such questions, we have a simpler \textit{QA-as-classification} pipeline. As shown in Figure \ref{fig:pipeline}, our overall model is thus a hybrid model containing the following elements: (i) a binary classifier for deciding whether the given question can be answered from a small fixed vocabulary or needs more complex reasoning, and (ii) a simpler \textit{QA-as-classification} model to answer questions of the former type, and (iii) a multi-staged model containing four components as described below to deal with complex reasoning questions.}

\subsection{Visual Elements Detection (VED)}
The data bearing elements of a plot are of 10 distinct classes: the title, the labels of the $x$ and $y$ axes, the tick labels or categories (e.g., countries) on the $x$ and $y$ axis, the data markers in the legend box, the legend names, and finally the bars and lines in the graph.
Following existing literature (\cite{ClicheRMY17}, \cite{DVQA}), we refer to these elements as the {\em visual elements} of the graph. 
The first task is to extract all these visual elements by drawing bounding boxes around them and classifying them into the appropriate class. To this end, we can either apply object detection models such as RCNN, Fast-RCNN \cite{FastRCNN}, YOLO \cite{YOLO}, SSD \cite{LiuAESRFB16}, \textit{etc.} or instance segmentation models such as Mask-RCNN \cite{MaskRCNN}. \revision{Upon comparing all methods, we found that Faster R-CNN \cite{DBLP:conf/nips/RenHGS15} model along with Feature Pyramid Network(FPN) \cite{DBLP:conf/cvpr/LinDGHHB17} performed the best and hence we used it as our VED module.} %that 
%We treat this as an object detection cum instance segmentation task and use the Mask-RCNN model \cite{MaskRCNN} to detect these elements in the graph. 
%Figure \ref{fig:pipeline} shows the expected output of this stage with detected visual elements.

% \begin{figure*}
% \begin{center}
% \includegraphics[width=0.75\textwidth]{images/pipeline.pdf}
% \end{center}
% \caption{Our proposed multi-staged modular pipeline for QA on scientific plots.}
% \label{fig:pipeline}
% \end{figure*}

%Pipeline Figure
\begin{figure*}
\begin{center}
\includegraphics[width=\textwidth]{images/pipeline.pdf}
\end{center}
\caption{Our proposed multi-staged modular pipeline for QA on scientific plots.}
\label{fig:pipeline}
\end{figure*}


\begin{table}[]
\footnotesize
\begin{center}
\begin{tabular}{|l|r|r|}
\hline
\textbf{Dataset Split} & \textbf{\#Images} & \textbf{\#QA pairs} \\ \hline
Train & 157,070 & 20,249,479 \\
Validation & 33,650 & 4,360,648 \\
Test & 33,657 & 4,342,514 \\ \hline
% Test-Novel & 48,444 & 1,593,138 \\ \hline
%Test-Hard & 1,000  & 3,000 \\ \hline
\textbf{Total} & 224,377 & 28,952,641 \\ \hline
\end{tabular}
\caption{\revision{PlotQA Dataset Statistics}}
\label{basic-stats-old}
\end{center}
\end{table}

\subsection{Object Character Recognition (OCR)}
Some of the visual elements such as title, legends, tick labels, \textit{etc.} contain numeric and textual data. For extracting this data from within these bounding boxes, we use a state-of-the-art OCR model \cite{Smith07}. We crop the detected visual element to its bounding box, convert the cropped textual image into gray-scale, resize and deskew it, and then pass it to an OCR module. 
Existing OCR modules perform well for machine-written English text, and indeed we found that a pre-trained OCR module\footnote{https://github.com/tesseract-ocr/tesseract} works well on our dataset. 

\subsection{Semi-Structured Information Extraction (SIE)}
The next stage of extracting the data into a semi-structured table is best explained with an example shown in Figure \ref{fig:pipeline}.
The desired output of SIE is shown in the table where the rows correspond to the ticks on the $x$-axis (1996, 1997, 1998, 1999), the columns correspond to the different elements listed in the legend (Brazil, Iceland, Kazakhstan, Thailand) and the $i$,$j$-th cell contains the value corresponding to the $x$-th tick and the $y$-th legend. 
The values of the $x$-tick labels and the legend names are available from the OCR module. 
The mapping of legend name to legend marker or color is done by associating a legend name to the marker or color whose bounding box is closest to the bounding box of the legend name.
% The mapping of legend name to legend marker or color remains to be done. 
% Specifically in the example, we have identified the bounding boxes of all the legend colors and the legend names but we need to now associate the right color to the right legend name. 
% We do this by associating a legend name to the marker or color whose bounding box is closest to the bounding box of the legend name. 
Similarly, we associate each tick label to the tick marker whose bounding box is closest to the bounding box of the tick label. 
For example, we associate the legend name Brazil to the color ``Dark Cyan'' and the tick label 1996 to the corresponding tick mark on the $x$-axis. 
With this we have the 4 row headers and 4 column headers, respectively.
To fill in the 16 values in the table, there are again two smaller steps. 
First we associate each of the 16 bounding boxes of the 16 bars to their corresponding $x$-ticks and legend names. 
A bar is associated with an $x$-tick label whose bounding box is closest to the bounding box of the bar. 
To associate a bar to a legend name, we find the dominant color in the bounding box of the bar and match it with a legend name corresponding to that color.
Second, we need to find the value represented by each bar. 
% For this, we use the tick marks on the $y$-axis as identified by the OCR module as a reference.  
We extract the height of the bar using bounding box information from the VED module and then search for the $y$-tick labels immediately above and below that height. 
We then interpolate the value of the bar based on the values of these bounding ticks.
With this we have the 16 values in the cells and thus have extracted all the information from the plot into a semi-structured table. The output of each stage  is discussed in the supplementary material.

\subsection{Table Question Answering (QA)}
The final stage of the pipeline is to answer questions on the semi-structured table. 
As this is similar to answering questions from the WikiTableQuestions dataset \cite{PasupatL15}, we adopt the same methodology as proposed in \cite{PasupatL15}. 
In this method, the table is converted to a knowledge graph and the question is converted to a set of candidate logical forms by applying compositional semantic parsing.
These logical forms are then ranked using a log-linear model and the highest ranking logical form is applied to the knowledge graph to get the answer. 
Note that with this approach the output is computed by a logical form that operates on the numerical data. 
This supports complex reasoning questions and also avoids the limitation of using a small answer vocabulary for multi-class classification as is done in existing work on VQA.
\textcolor{black}{There are recent neural approaches for answering questions over semi-structured tables such as \cite{NeelakantanLAMA16,HaugGG18}. Individually these models do not outperform the relatively simpler model of \cite{PasupatL15}, but as an ensemble they show a small improvement of only (1-2\%).}
% We must add that there are some recent neural approaches for answering questions over semi-structured tables such as \cite{NeelakantanLAMA16,HaugGG18}. However, these approaches take an ensemble of many models and outperform the relatively simpler model of \cite{PasupatL15} only by a small margin (1-2\%). 
% In the absence of an ensemble, these neural methods do not perform better than the method proposed in \cite{PasupatL15}. 
To the best of our knowledge, there is only one neural method \cite{krishnamurthy-etal-2017-neural} which outperforms \cite{PasupatL15}, but the code for this model is not available which makes it hard to reproduce the results.
% Hence we chose the model of \cite{PasupatL15}.% for answering questions on the generated tables.

\begin{figure*}
\begin{center}
\includegraphics[width=\textwidth]{images/new_pipeline_v1.pdf}
\end{center}
\caption{Our proposed model containing (i) a question classifier for deciding whether the question can be answered from a fixed vocabulary (orange) or needs more complex reasoning (green), (ii) \textit{QA-as-classification} model to answer questions of the former type, and (iii) \textit{multi-staged model} as a pipeline of perception and QA modules for answering complex questions.}
\label{fig:our_model}
\end{figure*} 

% Hybrid Model Figure
% \begin{figure*}
% \begin{center}
% \includegraphics[width=0.8\textwidth]{images/new_pipeline_v1.pdf}
% \end{center}
% \caption{Our proposed model containing (i) a question classifier for deciding whether the question can be answered from a fixed vocabulary (orange) or needs more complex reasoning (green), (ii) \textit{QA-as-classification} model to answer questions of the former type, and (iii) \textit{multi-staged model} as a pipeline of perception and QA modules for answering complex questions.}
% \label{fig:pipeline}
% \end{figure*} 

%-------------------------------------------------------------------------

\section{Experiments}
\label{sec:experiments}
% In this section we detail the data splits, baseline models, hyperparameter tuning and evaluation metrics. 
%\subsection{Setup}

\subsection{Train-Valid-Test Splits} 
By using different combinations of $841$ indicator variables and $160$ entities (years, countries, \textit{etc}), we created a total of $224,377$ plots. 
Depending on the context and type of the plot, we instantiated the $74$ templates to create meaningful \{question, answer\} pairs for each of the plots. 
%The average number of questions per plot is 129.
We created train (70\%), valid (15\%) and test (15\%) splits (Table \ref{basic-stats-old}). 
% The dataset, crowd-sourced questions and the model will be publicly available upon acceptance of this paper.
% In addition we will also release the 7,000 questions that we collected by crowd-sourcing to generate the question templates.
\PN{The dataset and the crowd-sourced questions can be downloaded from the link: \url{bit.ly/PlotQA}}

\subsection{Models Compared}
We compare the performance of the following models:\\
\noindent - \textit{\textbf{IMG-only}}: This is a simple baseline where we just pass the image through a VGG19 and use the embedding of the image to predict the answer from a fixed vocabulary.\\
\noindent - \textit{\textbf{QUES-only}}: This is a simple baseline where we just pass the question through a LSTM and use the embedding of the question to predict the answer from a fixed vocabulary.\\ 
\noindent - \textit{\textbf{SAN}} \cite{SAN}: This is an encoder-decoder model with a multi-layer stacked attention \cite{Attention} mechanism. It obtains a representation for the image using a deep CNN and a representation for the query using LSTM. It then uses the query representation to locate relevant regions in the image and uses this to pick an answer from a fixed vocabulary.\\
\noindent - \textit{\textbf{SANDY}} \cite{DVQA}: This is the best performing model on the DVQA dataset and is a variant of SAN. Unfortunately, the code for this model is not available and the description in the paper was not detailed enough for us to reimplement it.\footnote{We have contacted the authors and while they are helpful in sharing various details, they do not have access to the original code now.} Hence, we report the numbers for this model only on DVQA (from the original paper).\\
\MK{\noindent - \textit{\textbf{LoRRA}} \cite{singh2019TowardsVM}: This is the recently proposed model on the TextVQA dataset. It concatenates the image features extracted from pre-trained ResNet-152 \cite{DBLP:conf/cvpr/HeZRS16} model with the region based features extracted from Faster-RCNN \cite{Detectron2018} model. It then reads the text present in the image using a pre-trained OCR module and incorporates an attention mechanism to reason about the image and the text. Finally, it does multi-class classification where the answer either comes from a fixed vocabulary or is copied from the text in the image.
}\\
\MK{\noindent - \textit{\textbf{BAN}} \cite{DBLP:conf/nips/KimJZ18}: This model exploits bilinear interactions between two groups of input channels, \textit{i.e.}, between every question word (GRU \cite{DBLP:conf/emnlp/ChoMGBBSB14} features) and every image region (pre-trained Faster-RCNN \cite{DBLP:conf/nips/RenHGS15} object features). It then uses low-rank bilinear pooling \cite{DBLP:conf/iclr/KimOLKHZ17} to extract the joint distribution for each pair of channels. BAN accumulates 8 such bilinear attention maps which are then fed to a two-layer perceptron classifier to get the final joint distribution over answers from a fixed vocabulary.}\\
%\delete{
%\noindent - \textit{\textbf{VOES}}: This is our model as described in section \ref{sec:proposed_model} which is specifically designed for questions which do not have answers from a fixed vocabulary.
%}
%\delete{
%\noindent - \textit{\textbf{VOES-Oracle}}: 
% This is our model where the first three stages of VOES are replaced by an Oracle which computes the ground truth table from which the plot was generated. The QA model then answers questions from this ground truth table. With this we can exclusively evaluate the performance of the WikiTableQA model.
%This is our model where the first three stages of VOES are replaced by an Oracle, \textit{i.e.}, the QA model answers questions on a table that has been generated using the ground truth annotations of the plot. With this we can evaluate the performance of the WikiTableQA model when it is not affected by the VED model's errors.
%}
\noindent - \textit{\textbf{Our Model}}:  This proposed model shown in Figure~\ref{fig:our_model} with two model paths.
The training data for the binary classification is generated by comparing the performance of the individual models: For a given question, the label is set to 1 if the performance of QA-as-classification model is better than the multi-stage pipeline, and 0 otherwise. 
       We use an LSTM to represent the input question and then perform binary classification on this representation.

%\delete{
%Given the complementary strengths of SAN-VQA and VOES, we train a hybrid model with a binary classifier which given a question decides whether to use the SAN or the VOES model. The data for training this binary classifier is generated by comparing the predictions of a trained SAN model and a trained VOES model on the training dataset. For a given question, the label is set to 1 (pick SAN) if the performance of SAN was better than that of VOES. We ignore questions where there is a tie. The classifier is a simple LSTM based model which computes a representation for the question using an LSTM and uses this representation to predict 1/0. At test time, we first pass the question through this model and depending on the output of this model use SAN or VOES.
%}

%\subsection{Hybrid Model}
%Although our proposed pipeline as described above, performs reasonably well on questions for which the answers are to be calculated, it has some shortcomings. It fails to answer Yes/No questions. Also positional information is lost in the tables. For example, to answer the question, ``How many bars are there on the $4^{th}$ tick from the right?'', the plot is required. These restrictions prevent the pipeline model to correctly answer the structural understanding and data retrieval questions. To overcome these shortcomings, we propose a hybrid model, which learns to either trust our pipeline's predictions or SAN's predictions.

%\subsubsection{Question Meta-classifier}
%Based on the question and the ground-truth answer, we first train a LSTM based binary question classifier, which learns whether to trust the pipeline's predictions or to trust the SAN's predictions. 
%The model has an embedding layer which takes each word of the question as input and projects it into a 50-dimensional vector space. Embedding layer is followed by an LSTM layer of 128 dimensions and a dense layer of 256 units, with ReLu \cite{Relu} activation function and a dropout\cite{dropout} of 0.5. We trained the meta-classifier with RMSProp as the optimizer with initial learning rate 0.001 and batch size of 128 for 10 epochs.
%Validation accuracy of the meta-classifier on our dataset was $87.3\%$ .

%\subsubsection{SAN + Pipeline}
%Given a question at test-time, the question meta classifier first decides which of the two parallel networks - SAN or Pipeline should answer the question. Based on the classifier's output, the corresponding network produces the inference.

\subsection{Training Details}

%\noindent\textbf{IMG-only:} The image classifier in this baseline uses the image features extracted from a pre-trained VGG19 network. The image features are then flattened before passing through through $2$ fully connected layers of $512$ and $128$ dimensions. The resulting feature is then fed to the output layer with $1000$ nodes, each corresponding to a answer in the fixed vocabulary. The model is trained for 3 epochs using Adam \cite{Adam} with an initial learning rate of $0.001$.%, with ReLu non-linearity in the dense layers and a batch size of 32 for 3 epochs.

%\noindent\textbf{QUES-only:} The question classifier in this baseline model contains a 128-dimensional word embedding layer followed by an LSTM layer with 100 hidden units. The output is fed to the output layer with $1000$ nodes, each corresponding to a answer in the fixed vocabulary. The model is trained using Adam with an initial learning rate of $0.001$, with a dropout of $0.2$ in the LSTM layer and a batch size of 32 for 3 epochs. 
 
\noindent\textbf{SAN:} We used an existing implementation of SAN\footnote{https://github.com/TingAnChien/san-vqa-tensorflow} for the initial baseline results. Image features are extracted from the last pooling layer of VGG19 network. Question features are the last hidden state of the LSTM. Both the LSTM hidden state and 512-d image feature vector at each location are transferred to a 1024-d vector by a fully connected layer, and added and passed through a non-linearity (tanh). The model was trained using Adam \cite{Adam} with an initial learning rate of $0.0003$ and a batch size of 128 for 25,000 iterations.

\noindent\textbf{Our model:}
The binary question classifier in the proposed model contains a 50-dimensional word embedding layer followed by an LSTM with 128 hidden units. The output of the LSTM is projected to 256 dimensions and this is then fed to the output layer. The model is trained for 10 epochs using RMSProp with an initial learning rate of 0.001. Accuracy on the validation set is $87.3\%$.
Of the 4 stages of the multi-stage pipeline, only two require training, \textit{viz.}, Visual Elements Detection (VED) and Table Question Answering (QA). 
As mentioned earlier, for VED we train a \revision{variant of Faster R-CNN \cite{DBLP:conf/cvpr/LinDGHHB17} with FPN} using the bounding box annotations available in PlotQA. 
We trained the model with a batch size of 32 for $200,000$ steps. We used RMSProp with an initial learning rate of $0.004$. 
For Table QA, we trained the model proposed in \cite{PasupatL15} using questions from our dataset and the corresponding \textcolor{black}{ground truth tables}. 
%Since this model is computationally expensive with a high training time, we could train it using only $400,000$ questions from our training set. 

%Training details of the meta-classifier are mentioned in section 4.3.1

\subsection{Evaluation Metric} We used accuracy as the evaluation metric. Specifically, for textual answers (such as India, CO2, etc.) the model's output was considered to be correct only if the predicted answer exactly matches the true answer. However, for numeric answers with floating point values, an exact match is a very strict metric 
% (for example, if the predicted answer is 10.5 and the true answer is 10 then in most cases it would be acceptable). 
We relax the measure to consider an answer to be correct as if it is within 5\% of the correct answer. %we used exact match accuracy as the evaluation metric \textcolor{black}{We applied 5\% threshold on the model generated floating point numbers:} This means the model generated answer $(\hat{y})$ is marked correct if it is within the 5\% range of the ground-truth answer $(y)$ \textit{i.e.,} the model generated floating point answer will be marked correct is it satisfies the following condition:
%$$(y - y*0.05) \leq \hat{y} \leq (y + y*0.05)$$\\
%We did exact match for textual answers because:
%\begin{enumerate}
%\item if we keep levenshtein distance of 3 then all the misclassified Yes and No answers will be marked correct (as levenshtein distance between Yes and No is 3).

%\item if we keep levenshtein distance of 2 then all the answers like Male and Female will be marked correct.

%\item if we keep levenshtein distance of 1 then all the answers like CO2 and NO2 will be marked correct.
%\end{enumerate}

\subsection{Human Accuracy on PlotQA dataset}
\revision{To assess the difficulty of the PlotQA dataset, we report  human accuracy on a small subset of the Test split of the dataset. With the help of in-house annotators, we were able to evaluate $5,860$ questions grounded in $160$ images. Human accuracy on this subset is found to be $80.47\%$. We used the evaluation metric as defined in section 5.4. Most human errors were due to numerical precision as it is difficult to find the exact value from the plot even with a 5\% margin.}

%-------------------------------------------------------------------------

\section{Observations and Results}
% The results of our experiments are summarized in Tables \ref{dvqa-overall} to Table \ref{SIE_acc}. We discuss them below.


\noindent \textbf{1. Evaluating models on PlotQA dataset (Table \ref{plotqa-overall}):}
\textcolor{black}{The baselines IMG-only and QUES-only performed poorly with an accuracy of \revision{$4.84\%$ and $5.35\%$} respectively.}
%We then evaluate SAN, BAN, LoRRA and our model on this dataset. 
Existing models (SAN, BAN, LoRRA) perform poorly on this dataset. In particular, BAN and LoRRA have an abysmal accuracy of less than 1\%. This is not surprising given that both models are not designed to answer OOV questions. Further, the original VQA tasks for which BAN was proposed does not have any complex numerical reasoning questions as found in PlotQA. Similarly, LoRRA was designed only for text based answers and not for questions  requiring numeric reasoning. Note that we have used the original code \cite{singh2018pythia} released by the authors of these models. Given the specific focus and limited capabilities of these existing models it may even seem unfair to evaluate these models on our dataset but we still do so for the sake of completeness and to highlight the need for better models. Lastly, our model gives the best performance of 22.52\% on the PlotQA dataset.\\
\PN{\textbf{Ablation Study of proposed method}: Table \ref{ablation} presents the details of the ablation study of the proposed method for each question type (structural, data retrieval, reasoning) and each answer type (binary, fixed vocabulary, OOV). QA-as-classification performs very well on Yes/No questions and moderately well on Fixed vocab. questions with a good baseline aggregate accuracy of 7.76\%. It performs poorly on Open vocab. question, failing to answer almost all the 3,507,045 questions in this category. On the other hand, the QA-as-multi-stage pipeline fails to answer correctly any of the Yes/No questions, performs moderately well on Fixed vocab. questions, and answers correctly some of the hard Open vocab. questions. Our model combines the complementary strengths of QA-as-classification and QA-as-multi-stage pipeline achieving the highest accuracy of 22.52\%. In particular, the performance improves significantly for all Fixed Vocab. questions, while retaining the high accuracy of QA-as-classification on Yes/No questions and QA-as-multi-stage pipeline's performance on Open vocab. We acknowledge that the accuracy is significantly lower than human performance. This establishes that the dataset is challenging and raises open questions on models for visual reasoning.}

%\delete{There is a significant difference in the performance of VOES and VOES-Oracle across multiple question-answer types.
%This implies that the visual element detection in the VOES pipeline can be further improved. }
% We analyze this in the following.

% \noindent \textbf{2. Analysis of the VOES pipeline}
% We analyze the performance of the of visual element detection (VED) and OCR. \\
% - Table \ref{maskrcnn_ved} shows that the VED module performs reasonably well at an IOU of 0.5. 
% For higher IOUs of 0.8 and 0.9, the accuracy falls drastically. For instance, at IOU of 0.9, dotlines are detected with an accuracy of under 5\%.
% Clearly, such inaccuracies would lead to incorrect table generation and subsequent QA.
% This brings out an interesting difference between this task and other instance segmentation tasks where the margin of error is higher (where IOU of 0.5 is accepted).
% A small error in visual element detection as indicated by mAP scores of 80\% is considered negligible for VQA tasks, however for PlotQA small errors can cause significantly misaligned table generation and subsequent QA.
% Retraining VED model with a higher IOU of 0.75 only resulted in a small increase in accuracy (last row).
% \\
% - In Table \ref{ocr-acc} we evaluate the performance of the OCR module in standalone/oracle mode and pipeline mode.  In the oracle mode, we feed ground truth boxes to the OCR model whereas in the pipeline model we perform OCR on the output of the VED module. We observe only a small drop in performance, which indicates that the OCR module is robust to the reduction in VED module's accuracy at higher IOU as it does not depend on the class label or the exact position of bounding boxes. \\
% - In summary, a highly accurate VED for structured images (for elements such as bars, lines, dotlines, axis, tick labels) is an open challenge to improve reasoning over plots. 
% - In summary, we have high accuracies of the VED and OCR modules, and the VOES-Oracle model matches the accuracy on WikiTableQA \cite{PasupatL15}, but still the entire VOES pipeline has a low accuracy on Reasoning questions. 
% We attribute this to the low accuracy in generating structured tables from the VED and OCR outputs. 
% A small error in visual element detection as indicated by mAP scores of 80\% is considered negligible for VQA tasks, however for PlotQA small errors can cause significantly misaligned table generation and subsequent QA.
% We illustrate this with an example in the Supplementary Material.
% Highly accurate VED for structured images is an open challenge to improve reasoning over plots. 
% An alternative direction of future work is to avoid the intermediate table generation and perform QA with deep neural models on the output of VED and OCR modules.

\noindent \textbf{2. Analysis of the pipeline}
We analyze the performance of VED, OCR and SIE modules in the pipeline. \\
\noindent \textbf{VED:} Table \ref{maskrcnn_ved} shows that the VED module performs reasonably well at an Intersection Over Union (IOU) of 0.5.
For higher IOUs of 0.75 and 0.9, the accuracy falls drastically. For instance, at IOU of 0.9, dotlines are detected with an accuracy of under 20\%.
Clearly, such inaccuracies would lead to incorrect table generation and subsequent QA.
This brings out an interesting difference between this task and other instance segmentation tasks where the margin of error is higher (where IOU of 0.5 is accepted).
A small error in visual element detection as indicated by mAP scores of 75\% is considered negligible for VQA tasks, however for PlotQA small errors can cause significantly misaligned table generation and subsequent QA.
\textcolor{black}{We illustrate this with an example given in Figure \ref{VED}. The predicted red box having an IOU of 0.58 estimates the bar size as 760 as opposed to ground truth of 680, significantly impacting downstream QA accuracy.} 
\begin{figure}[h!]
\begin{center}
\includegraphics[scale=0.3]{images/bar_only.png}
\caption{Ground-truth (cyan) and predicted (red) boxes.}
\label{VED}
\end{center}
\end{figure}
%Retraining VED model with a higher IOU of 0.75 only resulted in a small increase in accuracy (last row).
% \textcolor{black}{Thus, inverting the plot generation function in going from the plot image to the structured table is a difficult CV task.}
\\
\noindent \textbf{OCR:} We evaluate the OCR module in standalone/oracle mode and pipeline mode in Table \ref{ocr-acc}.  In the oracle mode, we feed ground truth boxes to the OCR model whereas in the pipeline model we perform OCR on the output of the VED module. We observe only a small drop in performance from 97.06\% (oracle) to 93.10\% (after VED), which indicates that the OCR module is robust to the reduction in VED module's accuracy at higher IOU as it does not depend on the class label or the exact position of bounding boxes. \\
\MK{ 
\noindent \textbf{SIE:} We now evaluate the performance of the SIE module. We consider each cell in the table to be a tuple of the form \{row header, column header, value \} (e.g., \{Poland, 1964, 10000 tractors). We consider all the tuples extracted by the SIE module with the tuples present in the ground truth table to compute the F1-score. Even though Table \ref{maskrcnn_ved} suggests that the VED model is very accurate with a mAP@0.5 of 96.43\%, we observe that the F1-score for table extraction is only 0.68. This indicates that many values are not being extracted accurately due to the kind of errors shown in Figure \ref{VED} where the bounding box has a high overlap with the true box. We thus need better plot VED modules which can predict tighter bounding boxes (higher mAP at IOU of 0.9) around the plot's visual and textual elements. Inaccurate VED module leads to erroneous tables which further affects the downstream QA accuracy.}\\
- In summary, a highly accurate VED for structured images %(for elements such as bars, lines, dotlines, axis, tick labels) 
is an open challenge to improve reasoning over plots. 
% - In summary, we have high accuracies of the VED and OCR modules, and the VOES-Oracle model matches the accuracy on WikiTableQA \cite{PasupatL15}, but still the entire VOES pipeline has a low accuracy on Reasoning questions. 
% We attribute this to the low accuracy in generating structured tables from the VED and OCR outputs. 
% A small error in visual element detection as indicated by mAP scores of 80\% is considered negligible for VQA tasks, however for PlotQA small errors can cause significantly misaligned table generation and subsequent QA.
% We illustrate this with an example in the Supplementary Material.
% Highly accurate VED for structured images is an open challenge to improve reasoning over plots. 
% An alternative direction of future work is to avoid the intermediate table generation and perform QA with deep neural models on the output of VED and OCR modules.

\noindent \textbf{3. Evaluating new models on the existing DVQA dataset (Table \ref{dvqa-overall}):} 
% We first compare the performance of various models on the existing DVQA dataset (see, Table \ref{dvqa-overall}). 
The proposed model performs better than the existing models (SAN and SANDY-OCR) establishing a new SOTA result on DVQA. 
The higher performance of the proposed hybrid model in comparison to SAN (in contrast to the PlotQA results) suggests that the extraction of the structured table is more accurate on the DVQA dataset.
This is because of the limited variability in the axis and tick labels and shorter length (one word only) of labels. 
% On further analysis, we found that SAN and VOES complement each other wherein SAN typically performs better for \textit{Structural Understanding} questions and VOES performs better for \textit{Reasoning} questions. Hence, we chose hybrid model to combine the relative advantages of the two models. Indeed, 
%The hybrid model, SAN-VOES, improves on the individual models and establishes a new SOTA result on DVQA. 
% We present a detailed per question/answer type analysis later in the context of the PlotQA dataset. 
% E2E accuracy of different models on DVQA dataset.

\begin{table}[t]
\footnotesize
\begin{center}
\begin{tabular}{|l|r|r|}
\hline
\multicolumn{1}{|c|}{\textbf{Model}} & \multicolumn{1}{c|}{\begin{tabular}[c]{@{}c@{}}\textbf{DVQA} \textbf{(TEST)}\end{tabular}} & \multicolumn{1}{c|}{\begin{tabular}[c]{@{}c@{}}\textbf{DVQA} \textbf{(TEST-NOVEL)}\end{tabular}} \\ \hline
SAN & 32.1\% & 30.98\% \\
SANDY-OCR & 45.77\% & 45.81\% \\
% VOES & 48.62\% & 50.3\% \\ 
Our Model & \textbf{57.99\%} & \textbf{59.54\%} \\ \hline
\end{tabular}
\caption{Accuracy of different models on DVQA dataset.}
\label{dvqa-overall}
\end{center}
\end{table}

% % WACV-Round 1
% \begin{table}[t]
% \footnotesize
% \begin{center}
% \setlength\tabcolsep{2.5pt}
% \begin{tabular}{|l|l|rrr|}
% \hline
% \multicolumn{5}{|c}{\textbf{Accuracy (in \%)}} \vline \\
% \hline
% \parbox{1.3cm}{\textbf{Model}\\\textbf{(Agg. acc.)}} & \parbox{1.3cm}{\textbf{Q Type\textbackslash} \\\textbf{A type}} & \textbf{Structural} & \textbf{Data Retrieval} & \textbf{Reasoning}  \\ \hline
% % SAN
% \multirow{3}{*}{\parbox{1.3cm}{SAN \\ (46.54)}} 
% & Yes/No  & 94.01 & 95.35 & 66.02\\
% & Fixed vocab. & 77.3 & 32.06 & 29.27 \\
% & Open vocab. & NA & 0.9 & 3.4 \\ \hline
% % Our Model
% \multirow{3}{*}{\parbox{1.3cm}{VOES \\ (20.22)}}
% & Yes/No  & 0.00 & 0.00 & 0.00 \\
% & Fixed vocab. & 42.29 & 27.61 & 25.48  \\
% & Open vocab. & NA & 32 & 15.44  \\ \hline
% %VOES Oracle
% \multirow{3}{*}{\parbox{1.3cm}{VOES-Oracle \\ (32.55)}}
% & Yes/No  & 0.00 & 0.00 & 0.00   \\
% & Fixed vocab. & 42.55 & 41.79 & 47.74  \\
% & Open vocab. & NA & 83 & 36.99 \\ \hline
% % Hybrid Model
% \multirow{3}{*}{\parbox{1.3cm}{SAN-VOES \\ \textbf{(53.96)}}}
% & Yes/No  & \textbf{94.01} & \textbf{95.35} & \textbf{66.02} \\
% & Fixed vocab. & \textbf{81.66} & \textbf{40.6} & \textbf{35.74}  \\
% & Open vocab. & NA & \textbf{32.06} & \textbf{17.1} \\ \hline
% \end{tabular}
% \caption{Accuracy of different models on  PlotQA.} %Note that there are no open vocab. answers for \textit{Structural Understanding} question templates (see Table \ref{sample_qs_stats}).}
% \label{plotqa-overall}
% \end{center}
% \end{table}

\begin{table}[t]
\footnotesize
\begin{center}
\begin{tabular}{|l|c|c|c|c|c|c|}
\hline
\textbf{Models} & \textbf{IMG} & \textbf{QUES} & \textbf{BAN} & \textbf{LoRRA} & \textbf{SAN} & \textbf{\begin{tabular}[c]{@{}c@{}}Our\\ Model\end{tabular}} \\ \hline
\multicolumn{1}{|c|}{\begin{tabular}[c]{@{}c@{}}\textbf{Accuracy}\end{tabular}} & \multicolumn{1}{r|}{4.84} & \multicolumn{1}{r|}{5.35} & \multicolumn{1}{r|}{0.01} & \multicolumn{1}{r|}{0.02} & \multicolumn{1}{r|}{7.76} & \multicolumn{1}{r|}{\textbf{22.52}} \\ \hline
\end{tabular}
\caption{\revision{Accuracy (in \%) of different models on PlotQA.}}
\label{plotqa-overall}
\end{center}
\end{table}

%___________________________________________________________________________
% WACV-Round 2
\begin{table}[t]
\footnotesize
\begin{center}
\setlength\tabcolsep{2.5pt}
\begin{tabular}{|l|l|rrr|}
\hline
\multicolumn{5}{|c}{\textbf{Accuracy (in \%)}} \vline \\
\hline
\parbox{1.3cm}{\textbf{Model}\\\textbf{(Agg. acc.)}} & \parbox{1.3cm}{\textbf{Q Type\textbackslash} \\\textbf{A type}} & \textbf{Structural} & \textbf{Data Retrieval} & \textbf{Reasoning}  \\ \hline
% Human Accuracy
\multirow{3}{*}{\parbox{1.3cm}{\PN{Human\\ Baseline} \\ (80.47)}}
& Yes/No  & 99.77 & 100 & 76.51 \\
& Fixed vocab. & 99.29 & 83.31 & 59.97\\
& Open vocab. & NA & 87.58 & 58.01 \\ \hline
% SAN
\multirow{3}{*}{\parbox{1.3cm}{QA as classification \\ (7.76)}} 
& Yes/No  & 91.12 & 97.32  & 62.75 \\
& Fixed vocab. & 66.85 & 30.76 & 16.03  \\
& Open vocab. & NA & 0.00 & 0.00 \\ \hline
% VOES Oracle
% \multirow{3}{*}{\parbox{1.3cm}{\textcolor{red}{VOES}\\(Oracle) \\ ()}}
% & Yes/No  &  &  &  \\
% & Fixed vocab. & & & \\
% & Open vocab. & NA & &  \\ \hline
% VOES
\multirow{3}{*}{\parbox{1.3cm}{\PN{Multistage Pipeline} \\ (18.46)}}
& Yes/No  & 0.00 & 0.00 & 0.00 \\
& Fixed vocab. & 42.12 & 16.07 & 7.24  \\
& Open vocab. & NA & 57.39 & 14.95  \\ \hline
% Hybrid Model
\multirow{3}{*}{\parbox{1.3cm}{Our Model \\ \textbf{(22.52)}}}
& Yes/No  & \textbf{91.12} & \textbf{97.32} & \textbf{62.75} \\
& Fixed vocab. & \textbf{66.86} & \textbf{22.64} & \textbf{7.95}  \\
& Open vocab. & NA & \textbf{57.39} & \textbf{14.95} \\ \hline
\end{tabular}
\caption{\PN{Ablation study of proposed method on PlotQA. Note that there are no open vocab. answers for \textit{Structural Understanding} question templates (see Table \ref{sample_qs_stats}).}}
\label{ablation}
\end{center}
\end{table}

% % \noindent \textbf{4. Dissecting the performance of the 4 stages of VOES:} We now present the performance of each of the 4 stages of VOES. \\
% % - Table \ref{maskrcnn_ved} summarizes the performance of the VED stage for each of the 10 visual elements using an IOU of 0.5 and 0.8. We note that the performance is reasonably good for most classes but there is a clear scope for improvement in detecting lines, bars and dot lines. These 3 elements are crucial because all numeric values depend on the accurate detection of these elements.  \\



% % VED Module's Accuracy for different classes on TEST split of PlotQA dataset
\begin{table}[t]
% VED Modules Accuracy
%\begin{minipage}[t]{0.3\linewidth}
\footnotesize
\begin{center}
\begin{tabular}{|l|r|r|r|}
\hline 
\textbf{Class} & \textbf{AP@0.5} & \textbf{AP@0.75} & \textbf{AP@0.9} \\ \hline
Title & 100.00\% & 78.83\% & 0.22\% \\
Bar & 95.84\% & 94.30\% & 85.54\% \\
Line & 72.25\% & 62.04\% & 37.65\% \\
Dotline & 96.30\% & 95.14\% & 18.07\% \\
X-axis Label & 99.99\% & 99.99\% & 99.09\% \\
Y-axis Label & 99.90\% & 99.90\% & 99.46\% \\
X-tick Label & 99.92\% & 99.74\% & 96.04\% \\
Y-tick Label & 99.99\% & 99.97\% & 96.80\% \\
Legend Label & 99.99\% & 99.96\% & 93.68\% \\
Legend Preview & 99.95\% & 99.94\% & 96.30\% \\ \hline
\textbf{mAP} & \textbf{96.43\%} & \textbf{92.98\%} & \textbf{72.29\%}\\ \hline
% \textbf{mAP - IOU 0.75} & \textbf{92.69\%} & \textbf{79.58\%} & \textbf{42.88\%} \\ \hline
\end{tabular}
\caption{\MK{VED Module's Accuracy on PlotQA dataset}}
\label{maskrcnn_ved}
\end{center}
% \end{table}
%\end{minipage}
\end{table}
% % - Next, in Table \ref{ocr-acc} we evaluate the performance of the OCR module in standalone/oracle mode (first column) and pipeline mode (secodn column).  In the oracle mode, we feed ground truth boxes to the OCR model whereas in the pipeline model we perform OCR on the output of the VED module. We observe that there is not much drop in the performance of the OCR module even in the pipeline model which reiterated that the VED module is good at identifying bounding boxes containing textual elements.  \\
% % - Finally, in Table \ref{}, we evaluate the performance of the QA model in oracle and pipeline mode. Here, we see a significant drop in the performance in the pipeline mode. This indicates that the QA system have very good potential and if we can reduce the errors in the previous stage, then the final accuracy can be significantly better.

% OCR Module Accuracy on TEST split of PlotQA dataset
% \if 0
\begin{table}[t]
%\begin{minipage}[t]{0.3\linewidth}
\footnotesize
\begin{center}
\begin{tabular}{|l|r|r|}
\hline
% \textbf{Textual Class}
& \textbf{Oracle} & \textbf{After VED}\\ \hline
Title & 99.31\% & 94.6\% \\
X-axis Label & 99.94\% & 95.5\% \\
Y-axis Label & 98.43\% & 97.07\% \\
X-tick Label & 94.8\% & 91.38\% \\
Y-tick Label & 93.38\% & 88.07\% \\
Legend Label & 98.53\% & 91.99\% \\
\hline
\textbf{Total} & \textbf{97.06\%} &\textbf{ 93.10\%} \\
\hline
\end{tabular}
\caption{\PN{OCR Module Accuracy on the PlotQA dataset.}}
\label{ocr-acc}
\end{center}
%\end{minipage}
\end{table}
% \fi 
%----------------------------------------

% \begin{table}[]
% \footnotesize
% \begin{center}
% \begin{tabular}{|c|r|r|r|}
% \hline
% \textbf{\begin{tabular}[c]{@{}c@{}}IOU /\\ Metric\end{tabular}} & \multicolumn{1}{c|}{\textbf{@0.5}} & \multicolumn{1}{c|}{\textbf{@0.75}} & \multicolumn{1}{c|}{\textbf{@0.9}} \\ \hline
% \textbf{mAP} & 96.43\% & 92.98\% & 72.29\% \\ \hline
% \textbf{F1-score} & 0.68 & 0.679 & 0.57 \\ \hline
% \end{tabular}
% \caption{\revision{VED Module's Accuracy and SIE Module's F1 score on PlotQA dataset.}}
% \label{map_table_f1}
% \end{center}
% \end{table}
%------------------------------------------


% % \begin{itemize}
% %     %\item SAN, SANDY, VOES, SAN-VOES on DVQA
% %     %\item Zoom into SAN for 9 cells, indicates requirement to perform well on open vocab
% %     %\item Ques-only?, SAN, VOES, SAN-VOES on PlotQA (could not do SANDY)
% %     %\item Zoom into SAN-VOES 
% %     %\item Pipeline analysis of VOES - VED accuracy, OCR accuracy, Table accuracy, QA-Oracle Accuracy 
% %     \item VED with higher IOU 
% % \end{itemize}

% % \begin{itemize}
% %     \item VED is an open question to enable more accurate table extraction - tried other methods, tried training with higher IOU
% %     \item Can we do away with a tables - VED output directly to QA
% %     \item Some models which we could not reproduce or access - SANDY and latest deep neural methods for Table QA
% % \end{itemize}

% % Tables to include:
% % \begin{itemize}
% %     \item SAN, SANDY, VOES, SAN-VOES on DVQA - Table \ref{dvqa-overall}
% %     \item SAN, SAN-VOES zoomed into 9 cells each for PlotQA - Table \ref{hybrid-acc}
% %     \item Ques-only?, SAN, VOES, SAN-VOES on PlotQA - Table \ref{plotqa-overall}
% %     \item VED module accuracy(Table \ref{maskrcnn_ved}), OCR module accuracy (Table \ref{ocr-acc}), Table accuracy (Mean error percentage is reported in Table \ref{SIE_acc}), QA-oracle accuracy (Table \ref{sempre-acc})
% %     \item VED higher IOU training accuracy - Table \ref{high_iou_ved}
% % \end{itemize}

% % \textcolor{blue}{Table Generation Accuracy} : Given the ground-truth table and the predicted table, we first rearranged the rows and columns of the latter with respect to the former (by finding the levenshtein distance between the ground-truth and generated row/column headers). After finding this one to one correspondence between the cells of ground-truth and predicted tables, we compare the ground-truth cell value with the generated value. If the generated value is within $k\%$ of the ground-truth value, we consider the error to be zero for that cell. Likewise, we compute the error for all such cells and have reported the mean error percentage of this table generation process in Table \ref{SIE_acc}.
% %Mean error percentage between the ground-truth and the generated data values.
% \begin{table}[t]
% % SIE Module Accuracy
% %\begin{minipage}[t]{0.3\linewidth}
% \footnotesize
% \begin{center}
% \begin{tabular}{|c|r|}
% \hline
% \textbf{Threshold} &\textbf{Accuracy} \\ \hline
% 0.0 & 17\%\\ \hline
% 0.05\% & 16.78\% \\ \hline
% 0.1\% & 16.35\% \\ \hline
% \end{tabular}
% \caption{Mean error percentage between the ground-truth and the generated data values.}
% \label{SIE_acc}
% \end{center}
% %\end{minipage}
% \end{table}

% %-------------------------------------------------------------------------
% % \section{Discussion and Failure Cases}

% % \subsection{Stage-wise Analysis of our Pipeline}
% % \textcolor{blue}{
% % \begin{enumerate}
% %     \item Accurate numerical reasoning requires VED with high IOU, thus anchor-based and proposals-based methods do not work well
% %     \item Mask-RCNN performed the best with VED, but still struggles at higher IOUs
% %     \item High IOU VED is an open question
% %     \item Dealing with line and dotline plots: Indeed, solutions based on bounding-box do not work well for lines and dotline plots. This is a challenge when dealing with structured images and we do not currently have a solution for it.
% %     \item Structured data extraction given a set of VEDs is a challenge (small errors in VEDs can lead to large errors in structured tables). This implies better models are required that exploit the inter-dependence between VEs
% %     \item OCR module errors: Change in even a single character leads to drastically different tables.
% %     \item SIE module errors: A different representation of graphs needs to be explored. Tabular representation does not capture positional information (eg, which bar is to the left of the first bar?)
% %     \item Table QA on specific set of plot questions (eg, Yes/No questions) using deep NLP methods needs to be explored. 
% % \end{enumerate}
% % Why we choose Mask-RCNN over other anchor based object detection methods? and What are the errors that are getting propagated in the later stages?}
% % We report the per element average precision in Table x at different IOU thresholds. Specifically, if the predicted box has an IOU greater than a certain threshold with the groundtruth box, we classify it as a positive box. Before trying Mask-RCNN, we tried different region proposal and anchor based methods (Fast-RCNN, Faster-RCNN, YOLOv3, Retinanet, SSD) to detect the visual elements present in the plots. We report the mean average precision of all these methods in Table y. We ascribe the poor performance of these methods to the constraints imposed by the choice of proposals, anchor boxes in these detector networks.
% % We demonstrate this with an example shown in Figure \ref{fig:anchor_boxes}.
% % Note that there are several anchor boxes of the object detector that have some overlap with some of the bars in a horizontal bar plot. 
% % However, none have a satisfactory overlap (IOU $> 0.6$) exclusively with specific bars. 
% % This occurs due to the different aspect ratios of the bars. 
% % Clearly, the desired aspect ratios depend on the data: shorter and longer bars will need different anchor boxes, and so will small elements like ticks and legend previews.
% % To evaluate the possibility of data-driven configuration of anchor boxes, we plot in Figure \ref{fig:anchor_boxes_scatter} the aspect ratios of different visual elements of plots. 
% % We find that this spans a very large region, and further there are no clear clusters with a continuous range of representative aspect ratios. 
% % Supporting such a large number of anchor boxes through a network such as SSD \cite{LiuAESRFB16} will be computationally prohibitive.
% % Other limitation we found is that these networks provide a very large number of region proposals: on average $1,386$ proposals per plot. 
% % Given that the plot contains typically less than 50 visual elements, this is again a prohibitive number of region proposals.

% % Given this limitation of object detectors, we explored segmentation based object detectors (SegNet \cite{Segnet}, PixelLink \cite{pixelink}) which do not use pre-specified anchor boxes. These methods classify each image pixel into one of the 10 plot specific classes. Though the mean average precision of these methods is better than the anchor based methods, they fail to classify different instances of the same class. For example, in the Figure a, all pixels which are classified as bars are merged together.

% % This warrants the need to combine object detection methods with segmentation techniques. The current state of the art model, Mask-RCNN, is an instance segmentation model which does exactly the same and hence was a perfect choice for the VED module.

% % % Problem
% % The bounding boxes generated by Mask-RCNN are very sensitive to the IoU. Consider a bounding box which has \textit{``200000''} as its textual content. Now suppose the bounding box generated by Mask-RCNN covers only \textit{``20000''}. Even though the IoU will be greater than 0.5 the value extracted during the OCR stage will be significantly different that the actual value. This opens new research problems in the vision community to develop IoU insensitive models.

% %-----------------------CONCLUSION-----------------------------

\section{Conclusion}
We introduce the PlotQA dataset to reduce the gap between existing synthetic plot datasets and real-world plots and question templates.
Analysis of existing VQA models on PlotQA reveals that they perform poorly for Open Vocabulary questions. This is not surprising as these models were not designed to handle complex questions which require numeric reasoning and OOV answers. We propose a hybrid model with separate pipelines for handling (i) simpler questions which can be answered from a fixed vocabulary and (ii) complex questions with OOV answers. For OOV questions, we propose a pipelined approach that combines visual element detection and OCR with QA over tables. The proposed model gives state-of-the-art results on both the DVQA and PlotQA datasets. 
Further analysis of our pipeline reveals the need for more accurate visual element detection to improve reasoning over plots.



% %-----------------------TABLES-----------------------------


% Sample tables for DVQA and PlotQA
% \begin{figure*}
% \centering
% \begin{subfigure}[t]{.5\textwidth}
% \includegraphics[scale=0.3]{images/dvqa_table.png}
% \caption{DVQA table}
% \label{fig:dvqa_table}
% \end{subfigure}
% \begin{subfigure}[t]{.5\textwidth}
%   \includegraphics[scale=0.3]{images/reason_table.png}
%   \caption{PlotQA table}
%   \label{fig:plotqa_table}
% \end{subfigure}
% \caption{Sample tables generated after passing the DVQA and PlotQA images through our pipeline.}
% \label{fig:table_examples}
% \end{figure*}




% % % Table QA Accuracy for PlotQA dataset on Oracle Tables and Pipeline Generated Tables. Accuracy is calculated by keeping 5\% threshold.
% \begin{table}[t]
% % 5\% match
% \footnotesize
% \begin{center}
% \begin{tabular}{|l|rrr|}
% \hline
% & \multicolumn{3}{c}{\textbf{Oracle Tables}} \vline \\
% \hline
% Answer\textbackslash{}Template & \textbf{Structural} & \textbf{Data Retrieval} & \textbf{Reasoning} \\ \hline
% Yes/No  & 0.00\% & 0.00\% & 0.00\%   \\
% Fixed vocab. & 42.55\% & 41.79\% & 46.54\%  \\
% Open vocab. & NA & 83\% & 36.99\% \\ \hline
% \end{tabular}
% \caption{Table QA Accuracy for PlotQA dataset on Oracle Tables. Accuracy is calculated by keeping 5\% threshold.}
% \label{sempre-acc}
% \end{center}
% \end{table}


%-----------------------------------------------------------------------------
% mAP accuracy when MaskRCNN is trained at higher IOUs
% \begin{table}[t]
% \footnotesize
% \begin{center}
% \begin{tabular}{|l|r|r|}
% \hline
% \textbf{\begin{tabular}[c]{@{}l@{}}MaskRCNN \\ model\end{tabular}} & \multicolumn{1}{c|}{\textbf{\begin{tabular}[c]{@{}c@{}}Trained at \\ IOU 0.5\end{tabular}}} & \multicolumn{1}{c|}{\textbf{\begin{tabular}[c]{@{}c@{}}Trained at \\ IOU 0.75\end{tabular}}} \\ \hline
% mAP@0.5 & 94.21\% & 92.69\% \\ 
% mAP@0.8 & 80.44\% & 79.58\% \\
% mAP@0.9 & 38.86\% & 42.88\% \\ \hline
% \end{tabular}
% \caption{VED stage accuracy when trained at different IOUs}
% \label{high_iou_ved}
% \end{center}
% \end{table}



% %-------------------------------------------------------------------------
% \section{Future Work}

% %-------------------------------------------------------------------------
% \section{Conclusion}


{\small
\bibliographystyle{ieee}
\bibliography{main}
}

\newpage

% Dummy Table 1
% \begin{table}[]
% \begin{tabular}{|l|r|r|}
% \hline
% \textbf{Operator} & \multicolumn{1}{l|}{\textbf{Indoor}} & \multicolumn{1}{l|}{\textbf{Outdoor}} \\ \hline
% Idea & 3.55 & 3.45 \\
% Airtel & 3.62 & 3.38 \\
% Vodafone & 3.73 & 3.51 \\
% Reliance & 3.57 & 3.41 \\ \hline
% \end{tabular}
% \caption{Ground-truth Table}
% \label{tab:error_table_1_gt}
% \end{table}

% Dummy Table 2
% \begin{table}[]
% \begin{tabular}{|l|r|r|}
% \hline
% \textbf{Dperator} & \multicolumn{1}{l|}{\textbf{Indoo}} & \multicolumn{1}{l|}{\textbf{Outdoor}} \\ \hline
% Idea & 3.53 & 3.45 \\
% Airtel & 3.60 & 3.40 \\
% Vodafone & 4.0 & 3.35 \\
% Reliance & 3.62 & 3.40 \\ \hline
% \end{tabular}
% \caption{Predicted Table}
% \label{tab:error_table_1_generated}
% \end{table}

% Dummy Table 3
% %Error example 2nd one
% \begin{table}[]
% \begin{tabular}{|l|r|r|r|r|}
% \hline
% \textbf{Years} & \multicolumn{1}{l|}{\textbf{Kuwait}} & \multicolumn{1}{l|}{\textbf{Liberia}} & \multicolumn{1}{l|}{\textbf{Malaysia}} & \multicolumn{1}{l|}{\textbf{Peru}}\\ \hline
% 2007 & 98 & 150  & 166 & 424 \\
% 2008 & 98 & 150  & 145 & 424\\
% 2009 & 98 & 150  & 145 & 380\\
% 2010 & 98 & 150  & 145 & 380\\ 
% 2011 & 98 & 150.5  & 133 & 309\\
% 2012 & 98 & 150.5 & 133 & 293 \\ \hline
% \end{tabular}
% \caption{Ground-truth Table}
% \label{tab:error_table_2_gt}
% \end{table}

% Dummy Table 4
% \begin{table}[]
% \begin{tabular}{|l|r|r|r|r|}
% \hline
% \textbf{Years} & \multicolumn{1}{l|}{\textbf{Kuwait}} & \multicolumn{1}{l|}{\textbf{Liberia}} & \multicolumn{1}{l|}{\textbf{Malaysia}} & \multicolumn{1}{l|}{\textbf{Peru}}\\ \hline
% 2007 & 100- & 150  & 165 & 420 \\
% 200B & 100- &  - & 155 & 420\\
% -2009 & 100- &  - & 155 & 380\\
% 2010 & 100- &  - & 153 & 383\\ 
% 2011 & 100- & 155  & - & 310\\
% 2012 & 100- & 155 & - & 295 \\ \hline
% \end{tabular}
% \caption{Generated Table}
% \label{tab:error_table_2_generated}
% \end{table}

\section*{\textcolor{white}{Pass}}
\newpage
\section*{Appendices}
The supplementary material is organised in the following manner:
Section \ref{kg_creation} describes the methodology of knowledge graph creation from structured data.
In Section \ref{failure_cases}, we further analyse our proposed pipeline and discuss the errors in each stage.
In Section \ref{plot_samples}, we provide sample plots from the PlotQA dataset.
In Section \ref{templates}, we list all the 74 question templates that were formulated from the crowd sourced questions.
\section{Construction of knowledge graph from structured data} \label{kg_creation} 
Following \cite{PasupatL15}, we convert the semi-structured table into a knowledge graph which has two types of nodes \textit{viz.} row nodes and entity nodes. The rows of the table become row nodes, whereas the cells of each row become the entity nodes in the graph. Directed edges exist from the row nodes to the entity nodes of that column and the corresponding table column header act as edge-labels. An example of knowledge graph of the semi-structured table given in Figure \ref{fig:sie_table} is shown in Figure \ref{fig:flow}. For reasoning on the knowledge graph, we adopted the same methodology as given in \cite{PasupatL15}. The questions are converted to a set of candidate logical forms by applying compositional semantic parsing. Each of these logical forms is then ranked using a log-linear model and the highest ranking logical form is applied to the knowledge graph to get the final answer. 

% SIE Table
% \begin{table}[]
% \begin{tabular}{|l|r|r|r|r|}
% \hline
% \multicolumn{1}{|c|}{\textbf{Years}} & \multicolumn{1}{c|}{\textbf{Brazil}} & \multicolumn{1}{c|}{\textbf{Iceland}} & \multicolumn{1}{c|}{\textbf{Kazakhstan}} & \multicolumn{1}{c|}{\textbf{Thailand}} \\ \hline
% 1996 & 13.174405 & 7.895492 & 19.037112 & 8.821224 \\
% 1997 & 11.680978 & 7.642265 & 14.860660 & 9.322298 \\
% 1998 & 9.304022 & 5.221005 & 15.246865 & 8.174856 \\
% 1999 & 11.370439 & 5.453609 & 9.747697 & 9.685002 \\ \hline
% \end{tabular}
% \caption{}
% \label{tab:my-table}
% \end{table}

\begin{figure}[h]
\centering
\begin{subfigure}{.45\textwidth}
\centering
\includegraphics[scale=0.4]{supple_images/sie_table_1.png}
\caption{\PN{Semi-Structured Table}}
\label{fig:sie_table}
\end{subfigure}
\begin{subfigure}{.45\textwidth}
\centering
\includegraphics[scale=0.55]{supple_images/flow.pdf}
\caption{\PN{Knowledge Graph}}
\label{fig:flow}
\end{subfigure}
\caption{\PN{An example of the knowledge graph constructed from the semi-structured table.}}
\label{KG}
\end{figure}

\section{\PN{Some failure cases}}
\label{failure_cases}

In this section, we visualize the output of each stage in our multistage pipeline and show some interesting failure cases with the help of an example.

- \textbf{VED Stage:} Although the bounding boxes predicted by Faster R-CNN coupled with Feature Pyramid Network (FPN) fit reasonably well at an IOU of 0.5, it is not acceptable as the values extracted from these bounding boxes will lead to incorrect table generation and subsequent QA. Example: In Figure \ref{fig:error_analysis_ved_1}b, consider the bar representing the ``Indoor User Rating'' value for ``Vodafone''. The overlap between the ground-truth box (blue) and the predicted box (red) is higher than 0.5 but the values extracted from the detected box is $4.0$ as opposed to the actual value which is $3.73$. Another interesting failure case is shown in Figure \ref{fig:error_analysis_ved_2}b. There are multiple overlapping data points and the model is able to detect only one of the points. This leads to incomplete table generation as shown in Figure \ref{fig:error_analysis_sie_2}b where the values for Liberia for the years 2008, 2009 and 2010 could not be extracted. This small error might be acceptable for other VQA tasks but for PlotQA these small errors will escalate to multiple incorrect answers.

- \textbf{OCR stage:}
A slight misalignment in the bounding boxes predicted by VED module causes significant errors while extracting the text.
% As mentioned earlier, the IOU of $0.5$ is not a good metric for detecting visual elements in structured images as it leads to significant errors while extracting the text from the predicted bounding box.
% The bounding boxes predicted by Mask-RCNN fit reasonably well at an IOU of 0.5 but this is not acceptable always. 
Example: In Figure \ref{fig:error_analysis_ved_1}b, consider the box enclosing the legend label ``Indoor''. The rightmost edge of the predicted bounding box is drawn over the letter ``r'', which makes the OCR module incorrectly recognize the text as ``Indoo''. A similar error is made while performing OCR on the X-axis title which is read as ``Dperator'' instead of ``Operator''. \PN{Consider another example where there are misaligned bounding boxes on axes tick-labels as shown in Figure \ref{fig:error_analysis_ved_2}b. The values extracted are 200B, -2009 and -100 as opposed to the ground-truth values 2008, 2009 and 100. This slight error leads to incorrect column name in the subsequent generated tables (Figure \ref{fig:error_analysis_sie_1}b and Figure \ref{fig:error_analysis_sie_2}b) and incorrect answers to all the questions pertaining to these labels as shown in Table \ref{tab:error_analysis_table_1} and Tale \ref{tab:error_analysis_table_2}.
}

- \textbf{SIE stage:} Figure \ref{fig:error_analysis_sie_1}a shows the oracle table which is generated by using the ground-truth annotations and Figure \ref{fig:error_analysis_sie_1}b shows the table generated after passing the plot image through the different stages of our proposed multistage pipeline. It is evident from the generated table that the errors propagated from the VED and the OCR stage has lead to an incorrect table generation.

- \textbf{QA stage:} In Table \ref{tab:error_analysis_table_1} and Table \ref{tab:error_analysis_table_2} we compare the answer predictions made by different models with the ground-truth answer on randomly sampled questions. Note that, our proposed model combines the complementary strengths of both, QA-as-classification and QA as multistage pipeline, models.

%Error Image 1
\begin{figure*}
\centering
\begin{subfigure}{.45\textwidth}
\centering
\includegraphics[scale=0.4]{supple_images/0_vbar_categorical.png}
\caption{Input plot image}
\end{subfigure}
\hspace{0.5cm}
\begin{subfigure}{.45\textwidth}
\centering
\includegraphics[scale=0.4]{supple_images/error_analysis_1_generated_v2.png}
\caption{Few examples of the predicted bounding boxes}
\end{subfigure}
\caption{\PN{Errors made by the VED stage (highlighted in red). }}
\label{fig:error_analysis_ved_1}
\end{figure*}

\begin{figure*}
\centering
\begin{subfigure}{.45\textwidth}
\centering
\includegraphics[scale=0.45]{supple_images/error_analysis_table_1_gt.png}
\caption{Oracle table generated using ground-truth annotations}
\end{subfigure}
\hspace{0.5cm}
\begin{subfigure}{.45\textwidth}
\centering
\includegraphics[scale=0.45]{supple_images/error_analysis_table_1_generated.png}
\caption{Generated Semi-Structured Table}
\end{subfigure}
\caption{\PN{Errors made by the OCR and SIE stage (highlighted in red). Note that most of these errors have been propagated from the VED stage.}}
\label{fig:error_analysis_sie_1}
\end{figure*}


\begin{table*}[h!]
\centering
\begin{tabular}{|l|r|r|r|r|}
\hline
\multicolumn{1}{|c|}{\textbf{Question}} & \multicolumn{1}{c|}{\textbf{\begin{tabular}[c]{@{}c@{}}Ground \\ Truth\end{tabular}}} & \multicolumn{1}{c|}{\textbf{\begin{tabular}[c]{@{}c@{}}QA as\\ classification\end{tabular}}} & \multicolumn{1}{c|}{\textbf{\begin{tabular}[c]{@{}c@{}}Multistage \\ Pipeline\end{tabular}}} & \multicolumn{1}{c|}{\textbf{\begin{tabular}[c]{@{}c@{}}Our \\ Model\end{tabular}}} \\ \hline
\textbf{Q1.} What is the average indoor user rating per operator? & 3.62 & 500 & 3.54 & 3.54 \\ 
\textbf{Q2.} What is the total user rating for Vodafone in the graph? & 7.24 & 5 & 7.35 & 7.35 \\ 
\textbf{Q3.} What is the label or title of the X-axis? & Operator & Years & Dperator & Dperator \\ 
\textbf{Q4.} What is the indoor user rating of Airtel? & 3.62 & 0 & 3.60 & 3.60 \\ 
\textbf{Q5.} How many groups of bars are there? & 4 & 4 & 4 & 4 \\ 
\begin{tabular}[c]{@{}l@{}}\textbf{Q6.} What is the ratio of indoor user rating of Airtel to that\\ of the outdoor user rating of Vodafone?\end{tabular} & 1.03 & No & 3.35 & 3.35 \\ 
\begin{tabular}[c]{@{}l@{}}\textbf{Q7.} What is the difference between the highest and lowest\\ outdoor user rating?\end{tabular} & 0.13 & 1.5 & 0.10 & 0.10 \\ 
\begin{tabular}[c]{@{}l@{}}\textbf{Q8.} Does ``Indoor'' appear as one of the legend-labels in\\ the graph?\end{tabular} & Yes & Yes & 1.0 & Yes \\ 
\begin{tabular}[c]{@{}l@{}}\textbf{Q9.} For how many operators are the indoor user rating\\ greater than the average outdoor user rating taken \\ over all operators?\end{tabular} & 4 & 4 & 3.4 & 4 \\ 
\begin{tabular}[c]{@{}l@{}}\textbf{Q10.} Is the sum of outdoor user rating in Reliance and\\ Idea greater than the maximum outdoor user rating \\ across all operators?\end{tabular} & Yes & Yes & 6.85 & Yes \\ \hline
\end{tabular}
\caption{Answers predicted by different models on the sample questions.}
\label{tab:error_analysis_table_1}
\end{table*}


% Error Image 2
\begin{figure*}
\centering
\begin{subfigure}{.45\textwidth}
\centering
\includegraphics[scale=0.35]{supple_images/error_analysis_2.png}
\caption{Input plot image}
\end{subfigure}
\hspace{0.5cm}
\begin{subfigure}{.45\textwidth}
\centering
\includegraphics[scale=0.35]{supple_images/error_analysis_2_generated.png}
\caption{Few examples of the predicted bounding boxes}
\end{subfigure}
\caption{\PN{Errors made by the VED stage (highlighted in red).}}
\label{fig:error_analysis_ved_2}
\end{figure*}

\begin{figure*}
\centering
\begin{subfigure}{.45\textwidth}
\centering
\includegraphics[scale=0.4]{supple_images/error_analysis_table_2_gt.png}
\caption{Oracle table generated using ground-truth annotations}
\end{subfigure}
\hspace{0.5cm}
\begin{subfigure}{.45\textwidth}
\centering
\includegraphics[scale=0.4]{supple_images/error_analysis_table_2_generated.png}
\caption{Generated Semi-Structured Table}
\end{subfigure}
\caption{\PN{Errors made by the OCR and SIE stage (highlighted in red). Note that most of these errors have been propagated from the VED stage.}}
\label{fig:error_analysis_sie_2}
\end{figure*}

\begin{table*}[h!]
\centering
\begin{tabular}{|l|r|r|r|r|}
\hline
\multicolumn{1}{|c|}{\textbf{Question}} & \multicolumn{1}{c|}{\textbf{\begin{tabular}[c]{@{}c@{}}Ground \\ Truth\end{tabular}}} & \multicolumn{1}{c|}{\textbf{\begin{tabular}[c]{@{}c@{}}QA as\\ classification\end{tabular}}} & \multicolumn{1}{c|}{\textbf{\begin{tabular}[c]{@{}c@{}}Multistage \\ Pipeline\end{tabular}}} & \multicolumn{1}{c|}{\textbf{\begin{tabular}[c]{@{}c@{}}Our \\ Model\end{tabular}}} \\ \hline
\textbf{Q1.} How are the legend-labels stacked? & horizontal & horizontal & horizontal & horizontal \\
\begin{tabular}[c]{@{}l@{}}\textbf{Q2.} What is the time (in hours) required to\\ prepare and pay taxes in Kuwait in 2008?\end{tabular} & 98 & 100 & 100- & 100- \\
\begin{tabular}[c]{@{}l@{}}\textbf{Q3.} What is the difference between in time (in \\ hours) required to prepare and pay taxes in \\ Kuwait in 2009 and the time (in hours) required \\ to prepare and pay taxes in Liberia in 2008?\end{tabular} & -52 & -0.5 & 100- & 100- \\
\begin{tabular}[c]{@{}l@{}}\textbf{Q4.} What is the difference between the highest \\ and  the lowest time (in hours) required to \\ prepare and pay taxes in Peru?\end{tabular} & 131 & 500 & 125 & 125 \\
\begin{tabular}[c]{@{}l@{}}\textbf{Q5.} Is it the case that every year the sum of time \\ (in hours) required to prepare and pay taxes in \\ Peru and Kuwait is greater than the sum of the \\ time (in hours) required to prepare and pay \\ taxes in Liberia and Malaysia?\end{tabular} & Yes & Yes & 320 & Yes \\ \hline
\end{tabular}
\caption{Answers predicted by different models on the sample questions.}
\label{tab:error_analysis_table_2}
\end{table*}

% Sample Plots
\section{Samples from the PlotQA dataset}
\label{plot_samples}
\PN{Few examples of the \{plot, question, answer\} triplets from the PlotQA dataset are shown in Figure \ref{fig:supple_dataset_examples}. For each of the plots, most of the question templates discussed in section \ref{templates} are applicable but depending on the context of the plot, their language varies from it's surface form.}
\begin{figure*}
\centering
\begin{subfigure}{.45\textwidth}
\centering
\includegraphics[scale=0.3]{supple_images/181_line.png}
\caption*{\textbf{Q1}: What is the difference between price of gasoline in Switzerland and price of gasoline in Macedonia in 2008?\\
\textbf{A}: 0.15}
\end{subfigure}
\hspace{0.2cm}
\begin{subfigure}{.45\textwidth}
\centering
\includegraphics[scale=0.35]{supple_images/3_vbar_categorical.png}
\caption*{\textbf{Q2}: In how many years is the number of neonatal deaths in Cuba greater than 500? \\
\textbf{A}: 2
}
\end{subfigure}
\begin{subfigure}{.45\textwidth}
\centering
\includegraphics[scale=0.35]{supple_images/10_dot_line.png}
\caption*{
\textbf{Q3}:  What is the difference between the highest and the second highest amount of earnings from goods?\\
\textbf{A}: $0.9\mathrm{e}+5$}
\end{subfigure}
\hspace{0.2cm}
\begin{subfigure}{.45\textwidth}
\centering
\includegraphics[scale=0.35]{supple_images/149_hbar_categorical.png}
\caption*{\textbf{Q4}: What is the ratio of the sodium content in Sample 37 to that in Sample 33?\\
\textbf{A}: 1.086
}
\end{subfigure}
\begin{subfigure}{.45\textwidth}
\includegraphics[scale=0.35]{supple_images/1_hbar_categorical.png}
\caption*{
\textbf{Q5}: What is the average air quality index value for NHMC per hour?\\
\textbf{A}: 0.167
}
\end{subfigure}
\hspace{0.2cm}
\begin{subfigure}{.45\textwidth}
\includegraphics[scale=0.35]{supple_images/1304_line.png}
\caption*{\textbf{Q6}: Does the amount of banana production monotonically increase over the years?\\
\textbf{A}: Yes
}
\end{subfigure}
\end{figure*}

\begin{figure*}
\centering
\begin{subfigure}{.45\textwidth}
\centering
\includegraphics[scale=0.35]{supple_images/17_dot_line.png}
\caption*{
\textbf{Q7}: What is the difference between two consecutive major ticks on the Y-axis?\\
\textbf{A}: $2.000\mathrm{e}+5$
}
\end{subfigure}
\hspace{0.5cm}
\begin{subfigure}{.45\textwidth}
\centering
\includegraphics[scale=0.35]{supple_images/2_vbar_categorical.png}
\caption*{\textbf{Q8}: In how many cases, is the number of bars for a given year not equal to the number of legend labels?\\
\textbf{A}: 0
}
\end{subfigure}
\begin{subfigure}{.45\textwidth}
\centering
\includegraphics[scale=0.35]{supple_images/6_vbar_categorical.png}
\caption*{\textbf{Q9}: In how many countries, is the mineral rent (as \% of GDP) in 1970 greater than the average mineral rent (as \% of GDP) in 1970 taken over all countries?\\
\textbf{A}: 1
}
\end{subfigure}
\hspace{0.5cm}
\begin{subfigure}{.45\textwidth}
\centering
\includegraphics[scale=0.35]{supple_images/39_line.png}
\caption*{\textbf{Q10}: What is the total tuberculosis detection rate in Indonesia?\\
\textbf{A}: 101
}
\end{subfigure}
% Trying
\begin{subfigure}{.45\textwidth}
\centering
\includegraphics[scale=0.32]{supple_images/182_dot_line.png}
\caption*{\textbf{Q11}: Is it the case that in every year, the sum of the number of tourists in Costa Rica and Serbia greater than the number of tourists in Bhutan?\\
\textbf{A}: Yes
}
\end{subfigure}
\hspace{0.5cm}
\begin{subfigure}{.45\textwidth}
\centering
\includegraphics[scale=0.42]{supple_images/6_hbar_categorical.png}
\caption*{\textbf{Q12}: How many bars are there on the $2^{nd}$ tick from the top?\\
\textbf{A}: 2
}
\end{subfigure}
\caption{Sample \{plot, question, answer\} triplet present in the PlotQA dataset.}
\label{fig:supple_dataset_examples}
\end{figure*}

% Q-Templates
\section{Question Templates}
\label{templates}
In this section, we present the 74 question templates which we have used for the question generation. Note that, not all question templates are applicable to each and every type of plot. Also depending on the context of the plot, the question varies from the template's surface form.
\begin{enumerate}
    \item \textbf{Structural Understanding} :
    \begin{enumerate}[1.]
        \item Does the graph contain any zero values?
        \item Does the graph contain grids ?
        \item Where does the legend appear in the graph ?
        \item How many legend labels are there?
        \item How are the legend labels stacked?
        \item How many $<$plural form of X\_label$>$ are there in the graph?
        \item How many $<$figure-type$>$s are there?
        \item How many different colored $<$figure-type$>$s are there?
        \item How many groups of $<$figure-type$>$s are there?
        \item Are the number of bars on each tick equal to the number of legend labels?
        \item Are the number of bars in each group equal?
        \item How many bars are there on the $i^{th}$ tick from the left?
        \item How many bars are there on the $i^{th}$ tick from the right?
        \item How many bars are there on the $i^{th}$ tick from the top?
        \item How many bars are there on the $i^{th}$ tick from the bottom?
        \item Are all the bars in the graph horizontal?
        \item How many lines intersect with each other?
        \item Is the number of lines equal to the number of legend labels?
    \end{enumerate}
    
    \item \textbf{Data Retrieval} :
    \begin{enumerate}[1.]
        \item What does the $i^{th}$ bar from the left in each group represent?
        \item What does the $i^{th}$ bar from the right in each group represent?
        \item What does the $i^{th}$ bar from the top in each group represent?
        \item What does the $i^{th}$ bar from the bottom in each group represent?
        \item What is the label of the $j^{th}$ group of bars from the left?
        \item What is the label of the $j^{th}$ group of bars from the top?
        \item Does the $<$Y\_label$>$ of/in $<$legend-label$>$ monotonically increase over the $<$plural form of X\_label$>$ ?
        \item What is the difference between two consecutive major ticks on the Y-axis ?
        \item Are the values on the major ticks of Y-axis written in scientific E-notation ?
        \item What is the title of the graph ?
        \item Does $<$legend\_label$>$ appear as one of the legend labels in the graph ? 
        \item What is the label or title of the X-axis ?
        \item What is the label or title of the Y-axis ?
        \item In how many cases, is the number of $<$figure\_type$>$ for a given $<$X\_label$>$ not equal to the number of legend labels ?
        \item What is the $<$Y\_value$>$ in/of $<i^{th}$ X\_tick$>$ ?
        \item What is the $<$Y\_value$>$ of the $i^{th}$ $<$legend\_label$>$ in $<i^{th}$ X\_tick$>$ ?
        \item Does the $<$Y\_label$>$ monotonically increase over the $<$plural form of X\_label$>$ ?
        \item Is the $<$Y\_label$>$ of/in $<$legend\_label1$>$ strictly greater than the $<$Y\_label$>$ of/in $<$legend\_label2$>$ over the $<$plural form of X\_label$>$ ?
        \item Is the $<$Y\_label$>$ of/in $<$legend\_label1$>$ strictly less than the $<$Y\_label$>$ of/in $<$legend\_label2$>$ over the $<$plural form of X\_label$>$ ?
    \end{enumerate}
    
    \item \textbf{Reasoning} :
    \begin{enumerate}[1.]
        \item Across all $<$plural form of X\_label$>$, what is the maximum $<$Y\_label$>$ ?
        \item Across all $<$plural form of X\_label$>$, what is the minimum $<$Y\_label$>$ ?
        \item In which $<$X\_label$>$ was the $<$Y\_label$>$ maximum ?
        \item In which $<$X\_label$>$ was the $<$Y\_label$>$ minimum ?
        \item Across all $<$plural form of X\_label$>$, what is the maximum $<$Y\_label$>$ of/in $<$legend\_label$>$ ?
        \item Across all $<$plural form of X\_label$>$, what is the minimum $<$Y\_label$>$ of/in $<$legend\_label$>$ ?
        \item In which $<$singular form of X\_label$>$ was the $<$Y\_label$>$ of/in $<$legend\_label$>$ maximum ?
        \item In which $<$singular form of X\_label$>$ was the $<$Y\_label$>$ of/in $<$legend\_label$>$ minimum ?
        \item What is the sum of $<$title$>$ ?
        \item What is the difference between the $<$Y\_label$>$ in $<i^{th}$x\_tick$>$ and $<j^{th}$x\_tick$>$ ?
        \item What is the average $<$Y\_label$>$ per $<$singular form of X\_label$>$ ?
        \item What is the median $<$Y\_label$>$ ?
        \item What is the total $<$Y\_label$>$ of/in $<$legend\_label$>$ in the graph?
        \item What is the difference between the $<$Y\_label$>$ of/in $<$legend\_label$>$ in $<i^{th}$x\_tick$>$ and that in $<j^{th}$x\_tick$>$ ?
        \item What is the difference between the $<$Y\_label$>$ of/in $<$legend\_label1$>$ in $<i^{th}$x\_tick$>$ and the $<$Y\_label$>$ of/in $<$legend\_label2$>$ in $<j^{th}$x\_tick$>$ ?
        \item What is the average $<$Y\_label$>$ of/in $<$legend\_label$>$ per $<$singular form of X\_label$>$ ?
        \item In the year $<i^{th}$x\_tick$>$, what is the difference between the $<$Y\_label$>$ of/in $<$legend\_label1$>$ and $<$Y\_label$>$ of/in $<$legend\_label2$>$ ?
        \item What is the difference between the $<$Y\_label$>$ of/in $<$legend\_label1$>$ and $<$Y\_label$>$ of/in $<$legend\_label2$>$ in $<i^{th}$x\_tick$>$ ?
        \item In how many $<$plural form of X\_label$>$, is the $<$Y\_label$>$ greater than $<$N$>$ units ?
        \item Do a majority of the $<$plural form of X\_label$>$ between $<i^{th}$ x\_tick$>$ and $<j^{th}$ x\_tick> (inclusive/exclusive)  have $<$Y\_label$>$ greater than N $<$units$>$ ?
        \item What is the ratio of the $<$Y\_label$>$ in $<i^{th}$ x\_tick$>$ to that in $<j^{th}$ x\_tick$>$ ?
        \item Is the $<$Y\_label$>$ in $<i^{th}$ x\_tick$>$ less than that in $<j^{th}$ x\_tick$>$ ?
        \item In how many $<$plural form of X\_label$>$, is the $<$Y\_label$>$ of/in $<$legend\_label$>$ greater than $<$N$>$ $<$units$>$?
        \item What is the ratio of the $<$Y\_label$>$ of/in $<$legend\_label1$>$ in $<i^{th}$ x\_tick$>$ to that in $<j^{th}$x\_tick$>$?
        \item Is the $<$Y\_label$>$ of/in $<$legend\_label$>$ in  $<i^{th}$ x\_tick$>$ less than that in $<j^{th}$ x\_tick$>$ ?
        \item Is the difference between the $<$Y\_label$>$ in $<i^{th}$x\_tick$>$ and $<j^{th}$x\_tick$>$ greater than the difference between any two $<$plural form of X\_label$>$ ?
        \item What is the difference between the highest and the second highest $<$Y\_label$>$ ?
        \item Is the sum of the $<$Y\_label$>$ in $<i^{th}$x\_tick$>$ and $<(i+1)^{th}$x\_tick$>$ greater than the maximum $<$Y\_label$>$ across all $<$plural form of X\_label$>$ ?
        \item What is the difference between the highest and the lowest $<$Y\_label$>$ ?
        \item In how many $<$plural form of X\_label$>$, is the $<$Y\_label$>$ greater than the average $<$Y\_label$>$ taken over all $<$plural form of X\_label$>$ ?
        \item Is the difference between the $<$Y\_label$>$ of/in $<$legend\_label1$>$ in $<i^{th}$x\_tick$>$ and $<j^{th}$x\_tick$>$ greater than the difference between the $<$Y\_label$>$ of/in $<$legend\_label2$>$ in $<i^{th}$x\_tick$>$ and $<j^{th}$x\_tick$>$ ?
        \item What is the difference between the highest and the second highest $<$Y\_label$>$ of/in $<$legend\_label$>$ ?
        \item What is the difference between the highest and the lowest $<$Y\_label$>$ of/in $<$legend\_label$>$ ?
        \item In how many $<$plural form of X\_label$>$, is the $<$Y\_label$>$ of/in $<$legend\_label$>$ greater than the average $<$Y\_label$>$ of/in $<$legend\_label$>$ taken over all $<$plural form of X\_label$>$ ?
        \item Is it the case that in every $<$singular form of X\_label$>$, the sum of the $<$Y\_label$>$ of/in $<$legend\_label1$>$ and $<$legend\_label2$>$ is greater than the $<$Y\_label$>$ of/in $<$legend\_label3$>$ ?
        \item Is the sum of the $<$Y\_label$>$ of/in $<$legend\_label1$>$ in $<i^{th}$x\_tick$>$ and $<j^{th}$x\_tick$>$ greater than the maximum $<$Y\_label$>$ of/in $<$legend\_label2$>$ across all $<$plural form of X\_label$>$?
        \item Is it the case that in every $<$singular form of X\_label$>$, the sum of the $<$Y\_label$>$ of/in $<$legend\_label1$>$ and $<$legend\_label2$>$ is greater than the sum of $<$Y\_label$>$ of $<$legend\_label3$>$ and $<$Y\_label$>$ of $<$legend\_label4$>$ ?
    \end{enumerate}
    
\end{enumerate}

\end{document}
